\documentclass[11pt]{article}
\usepackage[margin=1in]{geometry}
\usepackage{amsmath,amssymb,amsthm,mathtools}
\usepackage[hidelinks,pdfusetitle]{hyperref}
\usepackage{microtype}
\usepackage{enumitem}
\usepackage{array}
\usepackage{booktabs}
\newcommand{\DOI}[1]{\mbox{\href{https://doi.org/#1}{\nolinkurl{#1}}}}
\newtheorem{theorem}{Theorem}[section]
\newtheorem{lemma}[theorem]{Lemma}
\newtheorem{corollary}[theorem]{Corollary}
\newtheorem{proposition}[theorem]{Proposition}
\newtheorem{observation}[theorem]{Observation}
\newtheorem{conjecture}[theorem]{Conjecture}
\theoremstyle{definition}
\newtheorem{definition}[theorem]{Definition}
\theoremstyle{remark}
\newtheorem{remark}[theorem]{Remark}
\DeclareMathOperator{\Tr}{Tr}
\DeclareMathOperator{\spec}{spec}
\DeclareMathOperator{\dist}{dist}
\DeclareMathOperator{\diag}{diag}
\DeclareMathOperator{\ad}{ad}
\newcommand{\E}{\mathbb{E}}
\newcommand{\R}{\mathbb{R}}
\newcommand{\C}{\mathbb{C}}
\newcommand{\N}{\mathbb{N}}
\newcommand{\Tpair}{T_{\mathrm{pair}}}
\newcommand{\Tcold}{T_{\mathrm{pair}}^{\mathrm{cold}}}
\newcommand{\Nodd}{N_{\mathrm{odd}}}
\newcommand{\Dop}{\Delta}
\newcommand{\Dir}{K_{\beta}^{\mathrm{dir}}}
\newcommand{\Kcov}{K_{\beta}^{\mathrm{cov}}}
\newcommand{\cold}{\mathrm{cold}}
\newcommand{\LF}{\mathrm{LF}}
\emergencystretch=1.5em
\title{Obstructions to lifting cold-phonon pairing bounds to the full Gibbs
state in the Holstein model}
\author{Davis Rens\\SuperMatics Inc.\\\href{https://orcid.org/0000-0001-5839-5197}{ORCID 0000-0001-5839-5197}}
\date{}
\begin{document}
\maketitle

\begin{abstract}
For the spinful Holstein model at large $U/\Omega$, a recently released
argument bounds the polaronic pair-decay temperature by
$2^{26}(t^{2}/U)e^{-7U/(8\Omega)}$, but only on a cold manifold that keeps the
odd-site occupation zero and the phonon occupation below
$M=\lfloor U/(64\Omega)\rfloor$ at every site. Removing that projection, so that
the same ceiling holds for the Gibbs state of the unprojected Hamiltonian, is
open, and this note does not close it. It records why five natural route families
fail, quantifies what any surviving route must control, and isolates the remaining
statement. First, five route families are obstructed at the level of specified
estimates: operator Schrieffer-Wolff and phonon-number Lyapunov bounds, Gaussian
domination of the Lang-Firsov influence functional, reflection positivity,
log-concavity of the phonon marginal, and cut Schrieffer-Wolff without a global
projection. Second, the cutoff-free directed pair-transfer kernel satisfies an
exact bound that carries the same $e^{-(U/\Omega)\coth(\beta\Omega/4)}$ exponent,
but the influence weight of an electron history does not factorize over disjoint
time segments: an explicit legal history with four bond charges has conditional
log weight exceeding the free directed maximum by at least $3/4$ for all
$U/\Omega\ge20$ at $\beta\Omega=1$, and a closed four-hop pair-return loop costs
action at most $3$ independently of $U/\Omega$. A facing-plateau lemma survives
with the weaker constant $3/22$ and does not add along a spatial skeleton. A
sixth route, written in the distortion coordinate, removes the cold projection at
the level of the two-site problem: it gives a non-asymptotic splitting bound, an
excursion bound at every temperature, and a volume-uniform many-body barrier
bound, but it stops at the same locality step, because local valley projectors
leak at order $t^{2}/U$ rather than exponentially.
Third, in the bare Lang-Firsov frame the population of the hot region
$N_{z}>M$ is at least $c_{2}(t/U)^{2}$ with $c_{2}=1/4096$, on the dimer and,
with volume-uniform constants, on every torus and every finite
vertex-transitive regular graph, at every temperature in a fixed window and for
all hopping below an explicit power of $\Omega^{3}/U^{2}$. A separate
second-order Duhamel calculation on the dimer gives the same conclusion with a
different constant. The channel weight of a first-order virtual transfer tends
to $1/4$ both in the bare frame and after the first Schrieffer-Wolff rotation
used in the cold-projected proof, so that rotation removes the electronic gap
from the denominator and leaves the virtual cloud intact. Fourth, the same wall
appears in the charge-order problem: the strong-coupling pseudospin theory of the
same model is a known easy-axis model, and the first full-Fock transform, an
explicit local frequency window and a fixed-history trace bound are supplied
here, but positive-temperature charge order remains open at a decorated-contour
lemma that the nearest rigorous theorem cannot supply. The open problem for the
ceiling is a normalized charged-skeleton influence bound, stated here precisely.
No pairing ceiling, and no bound on any critical temperature, is proved or
disproved.
\end{abstract}

\section{Question, scope, and prior work}

This note is the third in a series that separates what is known about pair
correlations in electron-phonon models from what is merely plausible. The first
\cite{barriers1} examined determinant zeros of single-spin propagators. The
second \cite{barriers2} examined integrated dynamic cumulants,
positive determinants, reflection positivity and the vanishing effective
attraction of bond coupling. The present note concerns a different and narrower
question: whether a bound that has been proved under a cold-phonon projection can
be transferred to the full Gibbs state of the same Hamiltonian.

\subsection{The released statement, and the step that is missing}

For the spinful Holstein model on the square torus, with polaronic energy
$U=F^{2}/K$, phonon frequency $\Omega$, hopping $t$, and $X=U/\Omega$, the
companion \cite{largeu} proves the following. Let
$M=\lfloor X/64\rfloor$. On the cold manifold that keeps every site doubly
occupied or empty and every site phonon number at most $M$, after a
first-order Schrieffer-Wolff rotation, the equal-time polaronic pair correlations
decay algebraically with an exponent that is positive, and above all temperatures
$\tau$ with
\begin{equation}
\tau>2^{26}\,\frac{t^{2}}{U}\,e^{-\frac{7}{8}X},
\label{eq:coldceiling}
\end{equation}
uniformly summable. In the notation recalled in Section~\ref{sec:model} this
reads $k_{\mathrm{B}}\Tcold\le 2^{26}(t^{2}/U)e^{-7U/(8\Omega)}$.

That bound is a bound on a projected quantity. The physically natural companion
statement, the same inequality for the pair-decay temperature $\Tpair$ of the
Gibbs state of the unprojected Hamiltonian, would be a genuine ceiling: at strong
coupling, that is $\Omega$ small compared with $U$ at fixed $t$ and $U$, it
narrows to zero as $\Omega\downarrow0$. Whether it holds is open. It is not
implied by \eqref{eq:coldceiling}, and the projected exponent is not a bound on
the discarded sector.

Two obstructions were already isolated in \cite{largeu} itself. The first is
thermal and is geometric in the volume: the Gibbs weight of the cold manifold of
the atomic reference vanishes like
$\exp(-|\Lambda|e^{-\beta\Omega(M+1)})$ at every fixed temperature, so restricting
the Gibbs trace to the cold manifold cannot reproduce a linear-in-$X$ exponent.
The second is a bond-local floor on two-site pair transfer. Neither says
anything about the size of the discarded sector inside the true Gibbs state. This
note records further obstructions, of a different character, together with
the local statements that survive them.

\subsection{Summary of the results}

The results fall into six groups, in the order in which the note presents them.

\begin{enumerate}[label=(\Alph*),leftmargin=2.2em]
\item \emph{Five route families are obstructed at the level of specified
estimates} (Section~\ref{sec:routes}). Operator-level Schrieffer-Wolff and
phonon-number Lyapunov bounds meet a Franck-Condon virtual cloud of weight
$\Theta((t/U)^{2})$ at phonon number of order $X$, and the full-space
Schrieffer-Wolff generator is not defined because of exact resonances. Gaussian
domination of the Lang-Firsov influence functional discards the only factor that
carries $\Omega$. Neither candidate reflection kernel is reflection positive on
the smallest two-dimensional torus. The integrated phonon marginal is not
log-concave. Cut Schrieffer-Wolff loses to an exponent counting argument on any
bond whose endpoint is hot. What survives is a single-site entropy bound that is
uniform in volume but of order $t/U$ per site, hence of no use on its own.
\item \emph{A distortion-coordinate route removes the cold projection on the dimer
and stops at locality} (Section~\ref{sec:valley}). Written in the physical
distortion coordinate, the dimer becomes a three-channel double well whose barrier
is $(U-t)^{2}/U$ and whose width is set by $U/\Omega$, so the splitting is
exponentially small in $U/\Omega$ without any projection. Three statements are
proved with no adiabatic approximation and no truncation of the electronic sector:
for the full $3\times3$ matrix operator the ground-doublet splitting obeys
$E_{1}-E_{0}\le Ue^{-c_{0}X}$ with
$c_{0}=\frac12c_{\mathrm{WKB}}(x,\frac18)\ge0.655327$ for every $X\ge31$; an
excursion bound holds at every temperature with an exact cancellation of the
$\beta\Omega$ factor; and the many-body Born-Oppenheimer barrier obeys a
volume-uniform tensorization inequality, so no collective distortion lowers it
below $U/2-4dt$ per site. The route fails one step later: local valley projectors
fail to commute with the boundary hopping, with a fixed-distortion three-site
leakage of $0.1408944\approx\sqrt{2}\,t^{2}/(2U)$, and inside a valley the neutral
fermion dynamics obeys $\beta t^{2}/U\gg1$ in the target strip.
\item \emph{The cutoff-free directed kernel is correct and insufficient}
(Section~\ref{sec:kernel}). The two-oscillator trace of two same-direction hops
has an exact closed form, and the integrated directed kernel obeys
\begin{equation}
\Dir:=2t^{2}\int_{0}^{\beta}e^{-Uu}F_{+}(u)\,du
\le\frac{2t^{2}}{U}\bigl(1-e^{-\beta U}\bigr)e^{-X\coth(\beta\Omega/4)} ,
\label{eq:dir}
\end{equation}
whose exponent is exactly the kind of narrowing one wants. The exact influence
action is recalled, including the displacement that the physical pair field
carries on its endpoints. But the weight of a chronological sequence of bond
charges does not factorize over disjoint time segments. An explicit legal
four-charge history exceeds the largest value any free directed block can take,
by an interval-enclosed margin of about $9.3$ at $U=40$, and by at least $3/4$ for
every $U/\Omega\ge20$ at $\beta\Omega=1$. A closed four-hop pair-return loop
costs total action at most $3$ regardless of $U/\Omega$. A facing-plateau lemma
does survive, with constant $3/22$, but its cost is local and does not add along a
skeleton.
\item \emph{The hot sector is not empty, at any temperature, and is
algebraically large} (Section~\ref{sec:vtail}). In the bare Lang-Firsov frame
the probability $\langle\mathbf{1}_{\{N_{z}>M\}}\rangle$ is at least
$c_{2}(t/U)^{2}$ with $c_{2}=1/4096$ on the dimer, on every torus with
volume-uniform constants, and on every finite vertex-transitive regular graph.
The proof is a local finite-rank commutator on the unbounded bosonic Fock space,
with no truncation of the Gibbs state. An independent second-order Duhamel
calculation on the dimer reaches the same conclusion with different constants,
$c_{2}^{\mathrm{Du}}=h/1024$ with $h=(1-e^{-a})^{2}$ at the lower endpoint $a$ of
the temperature window, and an explicit cubic remainder. The virtual channel weight of a first-order transfer tends to
$1/4$ in the bare frame and to $1/4$ after the first Schrieffer-Wolff rotation;
that rotation removes the electronic gap from the denominator and leaves the
Franck-Condon cloud untouched. What remains open is a nonperturbative lower bound
after that rotation.
\item \emph{Charge order stops at the same wall} (Section~\ref{sec:cdw}). In the
half-filled model the strong-coupling pseudospin theory is known to be an easy-axis
XXZ model with a strictly positive anisotropy at every finite frequency; that
identification is prior art, cited as such. What is added is a bounded
Lie-Schwinger generator on the unrestricted Fock space with volume-uniform
connected-word bounds, an explicit local frequency window
$U/8\le\omega\le U/4$ with $J_{xy}/J_{z}<1/16$, and a fixed-history Schatten
trace bound valid for $\beta D\ge4$. Positive-temperature charge-order long-range
order for the pure model remains open at a precisely stated decorated-contour
lemma, which the nearest rigorous bosonic contour theorem cannot supply because it
requires a positive direct repulsion.
\item \emph{The remaining statement is a charged-skeleton influence bound}
(Section~\ref{sec:openproblem}). Combining the failures above singles out one
missing lemma, which is stated explicitly: a deterministic skeleton for
long-lived pair transport, a bound on the signed sum of neutral histories
normalized by the true partition function, a bounded-overlap allocation of local
plateau costs, and a convergent remainder for short virtual histories. No
implication in either direction between that lemma and the lifted ceiling is
proved here.
\end{enumerate}

None of these statements is a physical counterexample, and none is evidence about
a material. Each concerns a precisely described class of estimates applied to a
precisely described model quantity. In particular, nothing here shows that the
lifted ceiling is false, and nothing here shows that a contour or polymer argument
cannot reach it. The cold-projected theorem of \cite{largeu} stands as released.

\section{Model, notation, and the lifted quantity}
\label{sec:model}

\subsection{The spinful Holstein model}

Fix an integer side $L\ge4$ and let $\Lambda=T_{L}=(\mathbb{Z}/L\mathbb{Z})^{2}$
be the square torus with nearest-neighbour edge set $E(\Lambda)$ and torus
distance $d(x,y)$. Each site carries one oscillator coordinate and a spinful
electronic level with four states. With $t\ge0$, $\Omega>0$,
$\alpha=\sqrt{U/(2\Omega)}$, $g=\sqrt{U\Omega/2}$ in the convention
$g(a_{x}^{\dagger}+a_{x})(n_{x}-1)$, and real $\mu$, the Hamiltonian is
\cite{holstein1959}
\begin{equation}
H_{\Lambda}=-t\sum_{e\in E(\Lambda)}\sum_{\sigma}B_{e,\sigma}
+\Omega\sum_{x\in\Lambda}a_{x}^{\dagger}a_{x}
-\sqrt{\frac{U\Omega}{2}}\sum_{x\in\Lambda}(n_{x}-1)\bigl(a_{x}^{\dagger}+a_{x}\bigr)
-\mu\sum_{x\in\Lambda}n_{x},
\label{eq:H}
\end{equation}
with $B_{e,\sigma}=c_{u\sigma}^{\dagger}c_{v\sigma}+c_{v\sigma}^{\dagger}c_{u\sigma}$
for $e=\{u,v\}$, $n_{x}=n_{x\uparrow}+n_{x\downarrow}$, and $U=F^{2}/K>0$ the
polaronic energy of the oscillator. There is no direct electron-electron
interaction. We write $\tau=k_{\mathrm{B}}T=\beta^{-1}$, $\hbar=1$, and the two
dimensionless ratios
\begin{equation}
x=\frac{t}{U},\qquad X=\frac{U}{\Omega} .
\label{eq:ratios}
\end{equation}
The equal-time on-site pair field is $\Dop_{x}=c_{x\downarrow}c_{x\uparrow}$, with
the operator order fixed once. The projector onto singly occupied sites is
$\mathbf{1}_{\{n_{x}=1\}}=1-(n_{x}-1)^{2}$, and
$\Nodd=\sum_{x}(1-(n_{x}-1)^{2})$.

Where a result holds for a general finite graph, for instance in
Section~\ref{sec:vtail}, we say so and take $\Lambda$ finite, simple,
vertex-transitive of degree $q$; the torus is the case $q=2d$ with $L\ge3$.
Nothing below uses two-dimensional reflection positivity except where that is the
object of study.

\subsection{Lang-Firsov frame and the cold manifold}

Let $D_{x}(\gamma)=\exp(\gamma(a_{x}^{\dagger}-a_{x}))$ be the Weyl displacement,
and let
\begin{equation}
U_{\LF}=\exp\Bigl(-\alpha\sum_{x\in\Lambda}(n_{x}-1)(a_{x}^{\dagger}-a_{x})\Bigr),
\qquad \alpha=\sqrt{\frac{X}{2}} ,
\label{eq:ulf}
\end{equation}
be the Lang-Firsov transformation \cite{langfirsov}. Since $U_{\LF}$ commutes with
$n_{x}$ and therefore with the charge,
\begin{equation}
U_{\LF}a_{x}U_{\LF}^{\dagger}=a_{x}+\alpha(n_{x}-1),\qquad
U_{\LF}c_{x\sigma}^{\dagger}U_{\LF}^{\dagger}=c_{x\sigma}^{\dagger}D_{x}(-\alpha),\qquad
U_{\LF}c_{x\sigma}U_{\LF}^{\dagger}=c_{x\sigma}D_{x}(+\alpha).
\label{eq:lfacts}
\end{equation}
With $\Omega\alpha^{2}=U/2$ and $(n_{x}-1)^{2}=1-n_{x,\mathrm{odd}}$ the linear
coupling cancels. Dropping the additive constant $-\tfrac U2|\Lambda|$,
\begin{equation}
H_{\LF}=U_{\LF}H_{\Lambda}U_{\LF}^{\dagger}+\tfrac U2|\Lambda|
=D_{0}+H_{\mathrm{ph}}+V_{\mathrm{mix}}+V_{\mathrm{diag}},
\label{eq:HLF}
\end{equation}
with $H_{\mathrm{ph}}=\Omega\sum_{x}n^{\mathrm{ph}}_{x}$, $D_{0}=\tfrac U2\Nodd-\mu N$,
and for each oriented edge, writing $A_{e}=\sum_{\sigma}c_{u\sigma}^{\dagger}c_{v\sigma}$,
\begin{equation}
V_{e}=-t\bigl(A_{e}D_{u}(-\alpha)D_{v}(\alpha)+\mathrm{h.c.}\bigr),
\qquad V_{\mathrm{mix}}=\sum_{e}\bigl(P_{e,0}V_{e}P_{e,2}+P_{e,2}V_{e}P_{e,0}\bigr),
\label{eq:Ve}
\end{equation}
where $P_{e,j}$ projects onto $n_{u,\mathrm{odd}}+n_{v,\mathrm{odd}}=j$ and
$V_{\mathrm{diag}}=\sum_{e}(V_{e}-V_{e,\mathrm{mix}})$. Writing
$A_{e,\mathrm{mix}}=P_{e,0}A_{e}P_{e,2}+P_{e,2}A_{e}P_{e,0}$ and
$A_{e,\mathrm{diag}}=P_{e,1}A_{e}P_{e,1}$, on the sixteen electronic states of a
bond $\|A_{e,\mathrm{mix}}\|_{\mathrm{op}}=\sqrt2$ and
$\|A_{e,\mathrm{diag}}\|_{\mathrm{op}}=1$.

For $X\ge1$ put $M=\lfloor X/64\rfloor$ and define the orthogonal projector
\begin{equation}
P_{\cold}=\prod_{z\in\Lambda}\Bigl((1-n_{z,\mathrm{odd}})\,
\mathbf{1}_{\{n^{\mathrm{ph}}_{z}\le M\}}\Bigr),
\qquad \mathcal{H}_{\cold}=\operatorname{ran}P_{\cold},
\label{eq:pcold}
\end{equation}
following \cite{largeu}. Let $S_{1}$ be the first-order Schrieffer-Wolff
generator of \cite{largeu}, which is the sum over bonds of
\begin{equation}
S_{1,e}=x\Bigl[\bigl(P_{e,0}A_{e}P_{e,2}-P_{e,2}A_{e}P_{e,0}\bigr)D_{u}(-\alpha)D_{v}(\alpha)-\mathrm{h.c.}\Bigr],
\qquad [S_{1},D_{0}]=-V_{\mathrm{mix}} .
\label{eq:S1}
\end{equation}
The cold-projected effective Hamiltonian is
$H_{\cold}=P_{\cold}e^{S_{1}}H_{\LF}e^{-S_{1}}P_{\cold}$ on the
finite-dimensional space $\mathcal{H}_{\cold}$, and its Gibbs state is denoted
$\langle\cdot\rangle_{\cold}$. The Schrieffer-Wolff transformation itself, and the
question of when such a rotation is a legitimate tool, is the subject of the
standard references \cite{bdl,dff,dls}.

\subsection{The two temperatures}

The model quantity that every result in this series concerns is a clustering
threshold for one two-point function, not a transition temperature.

\begin{definition}[clustering temperatures; \cite{largeu}]
\label{def:tpair}
Let $\Tpair$ be the infimum of the temperatures $\tau$ with the following
property: for every $\tau'>\tau$, every $L\ge4$, every real $\mu$ and every
$\Omega>0$ there are constants $A_{0}>0$ and $\xi>0$, independent of $L$, $\mu$
and $\Omega$, such that
$|\langle\Dop_{x}^{\dagger}\Dop_{y}\rangle|\le A_{0}e^{-d(x,y)/\xi}$ for all
distinct $x,y\in\Lambda$ and
$\sup_{L\ge4}\max_{x}\sum_{y}|\langle\Dop_{x}^{\dagger}\Dop_{y}\rangle|<\infty$,
the expectation being that of $H_{\Lambda}$. For fixed $U>0$, $t\ge0$ and
$\Omega>0$ with $X=U/\Omega\ge1$, let $\Tcold$ be the infimum of the temperatures
$\tau$ such that for every $\tau'>\tau$ there are constants $A_{0}>0$ and
$\xi>2$, independent of $L\ge4$ and $\mu\in\R$, with
$|\langle\Dop_{x}^{\dagger}\Dop_{y}\rangle_{\cold}|\le A_{0}(d(x,y)+1)^{-\xi}$
for all distinct $x,y$ and
$\sup_{L\ge4}\max_{x}\sum_{y}|\langle\Dop_{x}^{\dagger}\Dop_{y}\rangle_{\cold}|<\infty$.
\end{definition}

Both quantities are upper thresholds for uniform row summability of a single
two-point function. They say nothing about a phase transition and they are not
critical temperatures. The definition of $\Tpair$ is the one used in
\cite{largeu,strong}.

With this notation the released statement is the following, quoted from
\cite{largeu} in the form in which it is used below.

\begin{theorem}[cold-projected ceiling; \cite{largeu}]
\label{thm:cold}
Let $U>0$, $t\ge0$ with $x=t/U\le2^{-14}$, let $X=U/\Omega\ge1$ and
$M=\lfloor X/64\rfloor$. With $A(\beta)=8\beta U K_{4}(x)e^{-\frac78X}$ and the
algebraic exponent $\xi_{\cold}(\beta)$ of \cite{largeu}, the cold-state pair
correlations satisfy
$|\langle\Dop_{x_{0}}^{\dagger}\Dop_{y_{0}}\rangle_{\cold}|\le(d+1)^{-\xi_{\cold}(\beta)}$
for all distinct $x_{0},y_{0}$ with $d=d(x_{0},y_{0})\ge1$, on every $L\ge4$ and
every real $\mu$. If $\tau>2^{26}(t^{2}/U)e^{-\frac78X}$ then
$\xi_{\cold}(\beta)>2$ and the row sums are bounded uniformly in $L$ and $\mu$.
Consequently $k_{\mathrm{B}}\Tcold\le2^{26}(t^{2}/U)e^{-7U/(8\Omega)}$.
\end{theorem}

The lifted statement, which is open, is the same inequality for $\Tpair$ with no
projection. Writing it out:
\begin{equation}
\text{open:}\qquad k_{\mathrm{B}}\Tpair\ \le\ C\,\frac{t^{2}}{U}\,
e^{-\frac78 X}\quad\text{at fixed }t,U,\ \text{uniformly in }\Omega
\text{ and }\mu,
\label{eq:lift}
\end{equation}
with $C$ a numerical constant, or the corresponding statement that
$\Tpair/\Omega\to\infty$ is impossible at strong coupling. Nothing in this note
decides between them.

\subsection{Three different things that can go wrong}

It is worth separating three mechanisms, because they are commonly conflated and
because the obstruction appropriate to each is different.

\begin{enumerate}[label=(\roman*),leftmargin=2.2em]
\item \emph{Thermal.} The cold manifold has exponentially small Gibbs weight at
fixed temperature, by \cite[Prop.~3]{largeu}: for the atomic reference and
$\mu=0$ the projected-to-total ratio is
$[2(1-e^{-\beta\Omega(M+1)})/(2+2e^{-\beta U/2})]^{|\Lambda|}$, which behaves like
$\exp(-|\Lambda|e^{-\beta\Omega(M+1)})$ at fixed $\beta$. Any argument that
replaces the Gibbs trace by its restriction to the cold manifold therefore loses
an amount that is not recoverable site by site.
\item \emph{Quantum virtual.} Even at $t\to0^{+}$ and at any temperature, virtual
hopping populates the hot region $N_{z}>M$ with probability of order $(t/U)^{2}$
per bond, because the displaced oscillators of a virtual pair are centered at
phonon number of order $X$, far above $M=\lfloor X/64\rfloor$. This is
algebraically small but not exponentially small in $X$, so no expansion whose
activities must carry $e^{-cX}$ can be built on the bare frame.
Section~\ref{sec:vtail} proves this in the bare frame, and shows that the first
Schrieffer-Wolff rotation of \eqref{eq:S1} does not remove it.
\item \emph{Neutral history.} Once the electronic history is summed rather than
upper bounded, the weight of a chronology does not factorize over disjoint time
segments, and neutral excursions can enhance rather than suppress a charged
transfer. Section~\ref{sec:kernel} makes this quantitative. A cluster expansion
on the exact influence functional must resum these neutral histories against the
true partition function instead of multiplying free kernels.
\end{enumerate}

The five route families of Section~\ref{sec:routes} fail for reasons (i) and (ii)
chiefly, with (iii) deciding the fate of the one route that survives them, the
cutoff-free kernel route of Section~\ref{sec:kernel}.

Section~\ref{sec:routes} records the route-level obstructions,
Section~\ref{sec:valley} the distortion-coordinate route and its locality step,
Section~\ref{sec:kernel} the cutoff-free kernel and its conditional-factorization
failure, Section~\ref{sec:vtail} the virtual-tail lower bounds and the
Schrieffer-Wolff channel calculation, Section~\ref{sec:cdw} the charge-order
analogue, and Section~\ref{sec:openproblem} states the remaining open problem,
together with what is and is not claimed for it.

\section{Five route families and their obstructions}
\label{sec:routes}

\subsection{What is asserted in this section, and what is not}

Each of the five subsections below describes a family of estimates that has been
proposed, in one form or another, as a way to pass from Theorem~\ref{thm:cold} to
\eqref{eq:lift}. What is recorded is the specific inequality within that family
which fails, and the scale at which it fails. Three kinds of statement appear, and
they are labelled at each use:
\begin{itemize}[leftmargin=1.4em]
\item an \emph{exact identity or inequality}, proved in one line of algebra from
the definitions of Section~\ref{sec:model};
\item a \emph{finite witness}, meaning a computation on the specified two-, three-
or four-site matrix with all electron sectors and a stated phonon cutoff, whose
parameters are given in the text;
\item a \emph{scaling statement} about a class of estimates, for instance that a
bound of a given form can never reach a given size.
\end{itemize}
None of these is a theorem about the Gibbs state of $H_{\Lambda}$ in infinite
volume, and none excludes the existence of some other estimate of the same name;
it excludes the specified construction. Table~\ref{tab:routes} collects the
outcome. A sixth route family, written in the distortion coordinate rather than
in the phonon-number frame, is treated separately in Section~\ref{sec:valley}; it
does not fail at an estimate of the kind listed here, and it is not in the table.

\begin{table}[t]
\centering
\small
\caption{The five route families, the estimate inside each that fails, and the
scale of the failure at the stated parameters. Witnesses are finite-matrix
computations except where marked exact; $X=U/\Omega$ and $x=t/U$.}
\label{tab:routes}
\begin{tabular}{@{}>{\raggedright\arraybackslash}p{2.5cm}>{\raggedright\arraybackslash}p{5.0cm}>{\raggedright\arraybackslash}p{7.0cm}@{}}
\toprule
Route & Estimate that fails & Scale of the failure \\
\midrule
(a) operator
Schrieffer-Wolff, phonon-number Lyapunov & exponential moment
$\langle e^{sN_{z}}\rangle\le K(s)$ uniformly in $X$ &
virtual cloud of weight $\Theta(x^{2})$ at $N_{z}\sim X$; exact resonance
$U+\Omega(m-n)=0$; one-sided rotation gives
$\langle e^{N/4}\rangle=1.7\times10^{14}$ at $X=128$ \\
(b) Lang-Firsov path integral, Gaussian domination & $0\le e^{-\Phi_{\beta}}\le1$
used to discard the influence action &
the discarded action is the whole $\Omega$ dependence: exact attractive Hubbard
pair splitting $5.612$ at $U=40$, $t=8$, against a Holstein tunnel gap that is
already below double-precision resolution \\
(c) reflection positivity in the
phonon coordinate & chessboard bound on displaced configurations &
neither the direct nor the particle-hole reflection kernel is positive on
$C_{4}$: $\lambda_{\min}=-2.03\times10^{6}$ and $-4.6\times10^{3}$ at $t=1$,
$\beta=2$ \\
(d) Brascamp-Lieb, log-concavity of the phonon marginal & convexity
$V_{\mathrm{eff}}''\ge1/U$ &
$V_{\mathrm{eff}}''(0)=1/U-1/(2\tau)=-0.475$ at $U=40$, $\tau=1$: exact \\
(e) cut Schrieffer-Wolff without the global projection & hot bonds carry the
unrotated hopping $2t$ &
Dyson gain $e^{2t\beta}$ against penalty $e^{-X/64}$ crosses at
$x=1/128$: heuristic norm comparison \\
\bottomrule
\end{tabular}
\end{table}

\subsection{Route (a): operator Schrieffer-Wolff and phonon-number Lyapunov
functions}\label{sub:routea}

The proposal is to work on the full Fock space, throw away the projector
$P_{\cold}$ of \eqref{eq:pcold}, and replace it by a positive weight that grows
with phonon number: find $\varphi:\N\to[0,\infty)$ increasing, with
$\sum_{n}e^{-\varphi(n)}<\infty$, such that the Gibbs weight of a history which
visits phonon number $n$ is bounded by $e^{-\varphi(n)}$ times its bare hopping
weight, uniformly in $X$. Equivalently one wants an exponential moment bound
$\langle e^{sN_{z}}\rangle\le K(s)$ with $K$ independent of $X$, together with a
high-temperature or low-temperature splitting of the resulting inequality. The
energy of the atomic oscillator, displaced by $\alpha^{2}=X/2$ per occupied
neighbour direction, is what would supply $\varphi$.

Three separate things go wrong.

\emph{Wrong-sign tilting.} The natural object in this programme is the tilted
partition function $\Tr e^{-\beta H+sN}$, and Golden-Thompson gives
\begin{equation}
\Tr\bigl(e^{-\beta H+sN}\bigr)\le\Tr\bigl(e^{-\beta H}e^{sN}\bigr),
\label{eq:gt}
\end{equation}
which bounds the tilted partition function \emph{by} the moment one wants, not
the other way round. Relation \eqref{eq:gt} is exact for all $\beta,s$; the
inequality cannot be reversed, and Bogoliubov's inequality controls
$\langle H_{\mathrm{int}}\rangle$, not $\langle e^{sN}\rangle$.

\emph{Virtual cloud.} In the bare Lang-Firsov frame the first-order virtual
transfer of a doublon across a bond reaches phonon number of order
$\alpha^{2}=X/2$ on each endpoint with the Franck-Condon weight
\begin{equation}
w_{M}=t^{2}\sum_{n_{z}>M,\,n_{w}\ge0}\frac{e^{-2\alpha^{2}}\alpha^{2(n_{z}+n_{w})}}
{n_{z}!\,n_{w}!\,\bigl[U+\Omega(n_{z}+n_{w})\bigr]^{2}}
\;\xrightarrow[\Omega/U\to0]{}\;\frac{t^{2}}{4U^{2}} .
\label{eq:wM}
\end{equation}
The limit in \eqref{eq:wM} is a Poisson concentration statement for total mean
$2\alpha^{2}=X$ concentrated at $X$, and it is an identity at leading order in
$\Omega/U$. Consequently the exact ground-state expectations of a two-site dimer
at $U=40$, as $X$ varies from $10$ to $60$ with $t=8$, are
\begin{equation}
\langle n_{\mathrm{rel}}\rangle_{\mathrm{GS}}\approx0.195,\ 0.395,\ 0.795,\ 1.195,
\qquad
\langle e^{n_{\mathrm{rel}}/4}\rangle_{\mathrm{GS}}\approx4.47,\ 2.0\times10^{5},\
9.5\times10^{14},\ 4.61\times10^{24},
\label{eq:moments}
\end{equation}
the last value taken at a phonon cutoff large enough that the maximum of the
distribution lies inside the truncated space. The first moment agrees with second
order perturbation theory, $\langle n_{\mathrm{rel}}\rangle=\frac{t^{2}}{2U^{2}}X$,
to three digits at every $X$; the exponential moment is the quantity that a
Lyapunov function would have to control, and it grows by twenty orders of
magnitude over the same range.

The two requirements on $\varphi$ clash. To dominate the cloud at $n\sim X$ one
needs $e^{-\varphi(X)}\gtrsim x^{2}$, that is $\varphi(X)\lesssim2\ln(1/x)$, a
bound of order $\ln X$ that does not grow with $X$; to be increasing and
summable one needs $\varphi$ unbounded in $n$; and the observed moment growth in
\eqref{eq:moments} shows that no exponential moment $\langle e^{sN_{z}}\rangle$
with $s>0$ is bounded uniformly in $X$. A Lyapunov function cannot do both jobs.

\emph{Resonance and one-sided rotation.} The first-order generator
\eqref{eq:S1} has denominators $U+\Omega(m-n)$, where $n$ is the pair-sector
phonon number and $m$ the odd-sector one. On the full Fock space these vanish
exactly whenever $n-m=X$, which for integer $X$ has infinitely many solutions
$(n,m)\in\N^{2}$; $S_{1}$ is then undefined and no full-space Schrieffer-Wolff rotation
exists without further input. Restricting $S_{1}$ to the cold manifold, as in
\eqref{eq:S1}, leaves the states $|1,1;m\rangle$ with $m>M$ unrotated, so
\eqref{eq:wM} is untouched at $T=0$. A one-sided rotation that acts with the cold
projector on one endpoint only is well defined, since $U-\Omega M>0$, and it does
reduce the zero-temperature hot weight to $O(x^{4})$. It does not survive
$\tau\sim\Omega$: at $X=128$, $M=2$, $x=1/128$ and $\beta\Omega=1$, the
thermally populated states $|2,0;M+1\rangle$ are unrotated and hop to $m\sim X$,
producing $\langle e^{n_{\mathrm{rel}}/4}\rangle=1.69\times10^{14}$ at a phonon
cutoff inside which the maximum of the distribution is resolved, that is three
times the value obtained when the cutoff clips the maximum. Both numbers are far
outside anything an activity bound could tolerate.

\subsection{Route (b): Gaussian domination of the Lang-Firsov influence
functional}\label{sub:routeb}

The proposal is to write the pair correlation as a Feynman-Kac worldline sum
with the influence action $\Phi_{\beta}[q]\ge0$ of
\eqref{eq:Phi} and to use Gaussian domination, that is
$0\le e^{-\Phi_{\beta}[q]}\le1$, to move the action into a comparison measure.
Two things are lost, and both are fatal.

\emph{All of $\Omega$.} Setting $e^{-\Phi_{\beta}}\equiv1$ reduces the worldline
weight to an instantaneous attractive Hubbard model, in which $U$ survives and
$\Omega$ does not appear. On the dimer with $U=40$, $t=8$, the attractive Hubbard
pair splitting is the exact $(\sqrt{U^{2}+16t^{2}}-U)/2=5.612$, while the Holstein
tunnel gap of the same dimer is of order $(t^{2}/U)e^{-cX}$ at $X=40$, below the
resolution of double precision; the comparison therefore discards exactly the
factor that produces \eqref{eq:coldceiling}.

\emph{Short loops.} Keeping $\Phi_{\beta}$, the two simplest virtual excursions
rescaled to the same duration have actions that differ by a factor four, and one
of them is small. A pair excursion $(2,0)\to(0,2)\to(2,0)$, that is
$\Delta n=\pm2$ on each endpoint, of duration $\delta=1/U$ at $U=40$, $\Omega=1$,
$\beta=1$ costs $\Phi_{\beta}^{(2\text{-site})}(1/U)=3.892$, and the same
excursion across a thermal time $\delta=\beta/2$ costs $39.187$, asymptotic to
$4U\delta$. A neutral excursion $(2,0)\to(1,1)\to(2,0)$, that is
$\Delta n=\pm1$, costs exactly a quarter of those values, $0.973$ and $9.797$,
asymptotic to $U\delta$. Either way the action of a short excursion is $O(1)$ in
$X$ rather than $O(X)$, so it is not suppressed at strong coupling on the scale
that $\Phi_{\beta}$ would have to supply, and the equal-time pair correlation in
the physical frame stays $\Theta(x^{2})$: at $X=40$ on the dimer it is $0.0104$ at
$\tau=\Omega$, $0.0107$ at $\tau=t^{2}/U$ and $0.0102$ at $\tau=10^{-4}$,
remaining of order $x^{2}=0.04$ rather than decaying like $e^{-cX}$.

The route's useful residue is the directed kernel of
Section~\ref{sec:kernel}: the single factor $F_{+}$, which separates a charged
transfer from a neutral excursion, does carry the exponent, but it carries it for
one isolated hop and not for a worldline.

\subsection{Route (c): reflection positivity in the phonon coordinate}\label{sub:routec}

Since \eqref{eq:H} is stoquastic in the position basis, so that
$Z_{\mathrm{el}}(x)>0$ for every configuration $x$, one would like a chessboard
estimate: reflection positivity of the static kernel
$M_{ij}=Z_{\mathrm{el}}(x_{i},\theta x_{j})$ under a spatial reflection $\theta$,
which would bound the weight of large displacements $|x_{z}|>x_{0}$ uniformly in
volume. Reflection positivity is a strong tool of exactly this shape in spin
systems \cite{dls}, and it is the mechanism behind the released two-dimensional
bounds.

It fails here, already on the smallest two-dimensional torus. On $C_{4}$, the
$2\times2$ square with $t=1$, $\beta=2$ and the static field $b_{z}$ applied
directly at $g=2$, the reflection kernels for the four available symmetries have
smallest eigenvalues
\begin{equation}
\begin{aligned}
&\text{direct, edge bisecting:} && -2.028\times10^{6},\\
&\text{particle-hole, edge bisecting:} && -4631.7,\\
&\text{site centred:} && -3.776\times10^{5},\ -4.022\times10^{7}.
\end{aligned}
\label{eq:rp}
\end{equation}
The particle-hole kernel is positive semidefinite on the dimer,
$\lambda_{\min}=+0.0409$, so the failure is genuinely two dimensional and not a
normalization artifact: on $C_{4}$ the two halves are joined by two parallel
edges, and the worldline trace across two shared bonds does not factor as
$\Tr(K_{i}K_{j}^{\dagger})$ with a definite sign. Reflection positivity is
therefore not available as the volume-uniform mechanism in the Holstein
coordinate, and the route has no replacement for it.

\subsection{Route (d): Brascamp-Lieb and log-concavity of the phonon marginal}\label{sub:routed}

A second volume-uniform mechanism would be a Brascamp-Lieb variance bound, for
instance $\operatorname{Var}(b_{z})\le U\tau$ for the local field
$b_{z}$ conjugate to the displacement, provided a convexity estimate
$\nabla^{2}V_{\mathrm{eff}}\ge1/U$ holds on the phonon marginal. Conditional on a
fixed electronic occupation path the phonon action is exactly quadratic: in the
energy-scaled coordinate $b_{z}=\sqrt{U\Omega/2}\,(a_{z}+a_{z}^{\dagger})$, whose
bare potential is $b_{z}^{2}/(2U)$, the field is a Gaussian centred at the shifted
mean $U(n_{z}-1)$ with the occupation-independent quantum variance
$\frac{U\Omega}{2}\coth(\beta\Omega/2)$, whose static limit at $\beta\Omega\ll1$
is $U\tau$. That variance controls fluctuations around the shifted mean; it
gives no information about the occupation $n_{z}$ and hence none about the phonon
number in the bare frame, which is the only quantity that enters $P_{\cold}$.

Integrating the electrons out destroys the convexity. On one static site at
$U=40$, $\tau=1$,
\begin{equation}
V_{\mathrm{eff}}(b)=\frac{b^{2}}{2U}-2\tau\ln\Bigl(2\cosh\frac{\beta b}{2}\Bigr),
\qquad V_{\mathrm{eff}}''(0)=\frac1U-\frac{1}{2\tau}=-0.475<0
\quad\text{for all }\tau<U/2,
\label{eq:veff}
\end{equation}
with symmetric wells at $b\approx\pm U$ and, at $\tilde\mu=0$,
\begin{equation}
\operatorname{Var}(b)=\frac{U}{\beta}+\frac{U^{2}}{1+e^{-\beta U/2}}
=40+\frac{1600}{1+e^{-20}}=1639.9999967>U\tau,
\label{eq:varb}
\end{equation}
which is larger than the static value $U\tau$ by $U^{2}/(1+e^{-\beta U/2})$: exact
algebra followed by one-site arithmetic. On the dimer at $t=1$, $U=40$, $\tau=1$
both directional second derivatives at the origin are negative, with closed forms
$\frac1U-\frac{\tanh(\beta t/2)}{t}=-0.437117$ in the antisymmetric direction and
$\frac1U-\frac{\beta}{2\cosh^{2}(\beta t/2)}=-0.368224$ in the symmetric one; at
$t=4$ the same two closed forms give $-0.216007$ and $-0.010325$. The two wells are the
double occupancy and the empty site, that is precisely the electronic sector the
cold manifold was designed to remove, and no convexity estimate survives in its
presence.

\subsection{Route (e): cut Schrieffer-Wolff without the global projection}\label{sub:routee}

The last family keeps the algebraic structure of the cold-projected proof and
replaces the global projector by an edge-local one: on a bond whose endpoints are
cold, apply $S_{1,e}$; on other bonds, leave the hopping alone, and control the
difference by a Boltzmann penalty for the hot endpoint.

Three scales defeat it. First, a hot bond keeps $\|V_{e}\|_{\mathrm{op}}=2t$
unrotated, so a Dyson expansion on such a bond gains $e^{2t\beta}=e^{2xX}$ against
the penalty $e^{-M}\approx e^{-X/64}$, and the two exponents cross where
\begin{equation}
2x>\tfrac1{64}\quad\Longleftrightarrow\quad x>\tfrac1{128}.
\label{eq:threshold}
\end{equation}
This is a heuristic comparison of norms rather than a proof of defeat. For every
$X$ the gain beats the penalty once $x>1/128$, so at moderate $x$ the naive bound
cannot control hot bonds; for $x\le1/128$ the penalty beats the gain, and this
comparison by itself obstructs nothing. Since Theorem~\ref{thm:cold} needs
$x\le2^{-14}$, far below the crossing point, the obstruction inside the
cold-projected regime is carried by the two mechanisms that follow, neither of
which is a comparison of norms against the Boltzmann penalty. Second, one
application of the cut generator leaves a
residual coupling between the cold sector and the hot pair sector of norm
$\Theta(x^{2}U)$: on the dimer the one-step residual is $0.4637$ while
$t^{2}/U=0.4$, so no $e^{-cX}$ factor appears at this order. Third, on a three-site
chain with $U=40$, $\Omega=10$, $t=4$, $M=1$, the two edge generators cancel the
cold-to-cold-odd block to within double-precision roundoff but generate the
cross-commutator
$[S_{01},V_{12}]+[S_{12},V_{01}]$ with norm $0.984$ into the odd sector and
$0.534$ directly into $n_{2}>2M$. Cross-commutators are unavoidable whenever hot
sites are left unrotated, and they carry no exponential smallness.

\subsection{What survives: a volume-uniform single-site bound}

One positive statement of this family survives, and it is worth stating
carefully, because it is the only mechanism in this section that is uniform in
the volume without a projection. It is an entropy-energy bound, and it is
insufficient for a reason of size rather than of principle.

\begin{lemma}[single-site hot weight in the bare Lang-Firsov frame]
\label{lem:relent}
Let $\Lambda$ be a finite graph with $z$ of degree $d$, let $H_{\LF}$ be
\eqref{eq:HLF} and let $\rho_{\LF}=e^{-\beta H_{\LF}}/\Tr e^{-\beta H_{\LF}}$.
Split $H_{\LF}=H^{(z)}+V_{z}$, where $H^{(z)}$ omits the bonds incident to $z$.
Then $\|V_{z}\|\le 2dt$ and, for every $\beta>0$, $M\ge0$ and real $\mu$,
\begin{equation}
\bigl\langle\mathbf{1}_{\{N_{z}>M\}}\bigr\rangle_{\LF}
\le\min\Bigl\{1,\ \frac{2\|V_{z}\|+\tau\log2}{\Omega(M+1)}\Bigr\}
\le\min\Bigl\{1,\ \frac{4dt+\tau\log2}{\Omega(M+1)}\Bigr\}.
\label{eq:relent-bound}
\end{equation}
\end{lemma}

\begin{proof}
For each edge $e=\{z,w\}$ and spin $\sigma$, the operator
$c_{z\sigma}^{\dagger}c_{w\sigma}D_{z}(-\alpha)D_{w}(\alpha)+\mathrm{h.c.}$ has
norm $1$ because the displacements are unitary, so $\|V_{e}\|\le2t$ and
$\|V_{z}\|\le2dt$. Write $V=V_{z}$, $H=H_{\LF}$, $H_{\sigma}=H^{(z)}$, and let
$\rho_{\sigma}$ be the Gibbs state of $H_{\sigma}$. Expanding
$D(\rho\|\rho_{\sigma})=\Tr\bigl(\rho(\ln\rho-\ln\rho_{\sigma})\bigr)$ and using
$H-H_{\sigma}=V$ gives the exact identity
\begin{equation}
D(\rho\|\rho_{\sigma})=\ln Z_{\sigma}-\ln Z_{H}-\beta\Tr(\rho V)\le2\beta\|V\|,
\end{equation}
where $\ln Z_{\sigma}-\ln Z_{H}\le\beta\|V\|$ follows from
$H_{\sigma}\ge H-\|V\|$ and monotonicity of $A\mapsto\Tr e^{-A}$, and $-\Tr(\rho V)\le\|V\|$. Put
$p=\langle\mathbf{1}_{\{N_{z}>M\}}\rangle_{\LF}$ and
$q=\langle\mathbf{1}_{\{N_{z}>M\}}\rangle_{\sigma}$. Because $H^{(z)}$ omits all
bonds incident to $z$, the oscillator at $z$ decouples as a free harmonic
oscillator under $\rho_{\sigma}$, so the geometric sum gives
$q=e^{-\beta\Omega(M+1)}$. The binary relative entropy satisfies
$d(p\|q)=-h_{2}(p)+p\ln(1/q)+(1-p)\ln\bigl(1/(1-q)\bigr)\ge p\ln(1/q)-\ln2
=p\beta\Omega(M+1)-\ln2$, since $h_{2}(p)\le\ln2$ and $\ln(1/(1-q))\ge0$. By
monotonicity of relative entropy under the two-outcome measurement
$\{\mathbf{1}_{\{N_{z}>M\}},\mathbf{1}_{\{N_{z}\le M\}}\}$ \cite{ohya},
$d(p\|q)\le D(\rho\|\rho_{\sigma})$. Combining gives
$p\beta\Omega(M+1)\le2\beta\|V\|+\ln2$, which is \eqref{eq:relent-bound}.
\end{proof}

The bound is uniform in $|\Lambda|$ and in $\mu$. Its two regimes should be
distinguished, because one of them is empty of information. When $M$ is small,
specifically when $\Omega(M+1)\le2\|V_{z}\|+\tau\log2$, the right side of
\eqref{eq:relent-bound} exceeds one, the minimum returns the trivial bound $p\le1$, and
no finite-$M$ numerical check can certify anything. A check at $M=2$, $X=128$ on
the dimer falls in that regime and must not be quoted as evidence. An informative
check needs the opposite comparison, and one exists at strong coupling: on the
two-mode Lang-Firsov dimer with $U=16$, $\Omega=1$, $t=1/4$, $M=1$, $\tau=1$ and a
phonon space cut at $N_{\mathrm{ph}}=18$, where $J_{0}=\lVert V_{0}\rVert_{\mathrm{op}}=2t=1/2$,
one finds $\langle\mathbf{1}_{\{N_{0}>1\}}\rangle_{\LF}=0.135477$ against the
strictly smaller-than-one bound $(2J_{0}+\tau\log2)/(\Omega(M+1))=0.846574$.
In general \eqref{eq:relent-bound} is informative only once $M$ exceeds
$(2\|V_{z}\|+\tau\log2)/\Omega-1$, and even there it is linear in $1/M$. Its
content for the lift is therefore nil, because it does not decay in $X$: at
$M\simeq X/64$ the right side of \eqref{eq:relent-bound} is $\Theta(64d\,x)$, and
$\Theta(\tau/U)$ at high temperature, independently of how large $X$ is. What a polymer or contour expansion on the
full Gibbs state needs is an activity for a hot bond that is exponentially small
in $X$, so that the hot sector can be removed from the long-distance asymptotics
at a cost that beats any volume factor. Lemma~\ref{lem:relent} supplies none. It
also gives only an upper bound on the probability and no information about the
influence of a hot region on a pair that propagates at distance $d\to\infty$,
which is the quantity that actually enters a ceiling. Section~\ref{sec:vtail}
shows that no argument of this type can be tightened into an exponential bound in
the bare frame, because the hot weight is bounded \emph{below} by
$\Theta(x^{2})$ at every temperature.

\section{The distortion-valley route: lifting the projection on the dimer and
stopping at locality}
\label{sec:valley}

The five routes of Section~\ref{sec:routes} all work inside the phonon-number
frame, where the cold projection $P_{\cold}$ is the price of a convergent
expansion. A sixth route abandons that frame. It treats the static distortion as
the coordinate, in which the two configurations of a bond are two wells of a
Born-Oppenheimer surface separated by an explicit barrier, and the tunnelling
amplitude between them is exponentially small in $U/\Omega$ already at the level
of the two-site operator. On the dimer the projection can then be dropped without
compensation. This section records what that route proves and where it stops. It
is the only route in this note that produces a bound of the released form without
the cold manifold, and it is included for that reason. Its stopping point turns
out to be a locality statement of the same shape as Section~\ref{sec:openproblem}.

\subsection{Distortion coordinates and the dimer surface}

Under $U=2g^{2}/\Omega$, $g=\sqrt{U\Omega/2}$, the unitary change of variable
$v_{x}=\pm\sqrt{U\Omega/2}\,(a_{x}+a_{x}^{\dagger})$ maps the spinful Hamiltonian
of Section~\ref{sec:model} to
\begin{equation}
H_{\Lambda}
=-\frac{U\Omega^{2}}{2}\sum_{x\in\Lambda}\partial_{v_{x}}^{2}
+\sum_{x\in\Lambda}\frac{v_{x}^{2}}{2U}+H_{\mathrm{el}}(v)
-\frac{\Omega|\Lambda|}{2},
\label{eq:distortion}
\end{equation}
with $[v_{x},\partial_{v_{x}}]=-1$ and $H_{\mathrm{el}}(v)$ the electronic
Hamiltonian at frozen distortion $v$. This is the standard Holstein change of
variables \cite{holstein1959,langfirsov} and is used here only as a rewritten
frame; every statement below is about \eqref{eq:distortion} with the phonons kept
in the algebra.

On the dimer, with $c=(v_{1}+v_{2})/2$ and $r=(v_{1}-v_{2})/2$, the
centre-of-mass coordinate decouples as a free oscillator and the relative
coordinate carries a three-channel electronic matrix. In the spin-triplet sector
the electronic energy vanishes and the potential is $r^{2}/U$. In the singlet
sector, in the diabatic basis $\{|2,0\rangle,|1,1\rangle,|0,2\rangle\}$, the
matrix has eigenvalues $\{-2\sqrt{r^{2}+t^{2}},\,0,\,+2\sqrt{r^{2}+t^{2}}\}$ and
the lowest sheet is
\begin{equation}
V_{-}(r)=\frac{r^{2}}{U}-2\sqrt{r^{2}+t^{2}} .
\label{eq:vminus}
\end{equation}
For $U>t$ it has symmetric minima at $r=\pm r_{0}=\pm\sqrt{U^{2}-t^{2}}$ of depth
$V_{\min}=-U-t^{2}/U$, against a single-polaron minimum of $-U/2-t^{2}/U$ for
$U>2t$ (and $-U/4-t$ for $t<U\le2t$), so the two-site singlet is the ground
sector and is separated from the single-polaron sector by $U/2$; the well has
frequency
$\Omega_{\mathrm{eff}}=\Omega\sqrt{1-(t/U)^{2}}$, and the barrier between them is
exactly $B_{\mathrm{dimer}}=(U-t)^{2}/U$. The distortion picture thus converts an
$X$-dependent phonon-number bookkeeping problem into a double-well splitting
problem with a barrier linear in $U$ and a frequency $\Omega$: the exponent that
appears naturally is $U/\Omega$ itself, which is what the cold projection was
patching.

\subsection{Non-asymptotic splitting and excursion bounds}

The following three statements hold for the full $3\times3$ matrix operator
\eqref{eq:vminus} on $L^{2}(\mathbb{R})\otimes\mathbb{C}^{3}$, with no adiabatic
approximation and no truncation of the electronic sector. They are proved with
matrix Agmon estimates, a constant-spinor trial state, a global quadratic lower
bound on $V_{-}-V_{\min}$, parity pigeonholing and a Dirichlet-Neumann min-max
argument. The double-well tunnelling framework is classical
\cite{simon1984,hs1984,harrell1980,martinez1989,kmsm1992}; what is added here is a
non-asymptotic form with explicit constants in the Holstein parameters.

\begin{theorem}[ground doublet]\label{thm:valleyground}
Let $x=t/U\in(0,1/4]$, $X=U/\Omega$, and
\begin{equation}
X_{0}(x)=\max\Bigl\{\frac{8}{(1-x)^{2}},\
\frac{4}{s_{*}(1/8)^{2}},\
\frac{64(1-x)^{2}}{c_{\mathrm{WKB}}(x,1/8)^{2}}\Bigr\}\le31 ,
\label{eq:X0}
\end{equation}
with $s_{*}$ the Agmon parameter and $c_{\mathrm{WKB}}$ the classical
instanton normal form. Then for all $X\ge X_{0}(x)$,
\begin{equation}
E_{1}-E_{0}\le U\exp\bigl(-c_{0}(x)X\bigr),\qquad
c_{0}(x)=\tfrac12\,c_{\mathrm{WKB}}\bigl(x,\tfrac18\bigr)\ge0.655327 .
\label{eq:valleyground}
\end{equation}
As $X\to\infty$ at fixed $x$ the sharp exponent
$c_{\mathrm{WKB}}(x,0)$ is attained up to $O(X^{-1/4})$.
\end{theorem}

\begin{theorem}[all sub-barrier doublets]\label{thm:valleydoublets}
Fix $\kappa\in[0,1)$ and $x=t/U\in(0,1/4]$, and let $k\ge0$ obey
$E_{2k}\le V_{\min}+\kappa B_{\mathrm{dimer}}$ and $k+1\le11X$. Then for
$X\ge X_{0}(x,\kappa)$,
\begin{equation}
E_{2k+1}-E_{2k}\le C(x,\kappa)\,U\,\exp\bigl(-c(x,\kappa)X\bigr),\qquad
c(x,\kappa)=\tfrac12\,c_{\mathrm{WKB}}\Bigl(x,\tfrac{1+\kappa}{2}\Bigr),
\label{eq:valleydoublets}
\end{equation}
with the prefactor
\begin{equation}
C(x,\kappa)=\frac{160}{e\,c_{\mathrm{WKB}}\bigl(x,\tfrac{1+\kappa}{2}\bigr)}\,
\frac{352}{7(1-\kappa)(1-x)^{2}}.
\label{eq:valleyC}
\end{equation}
\end{theorem}

The exponent in \eqref{eq:valleydoublets} is the same matrix Agmon estimate as in
the ground case; what needs care is the eigenvalue count that fixes the prefactor.
Half-line Neumann bracketing at the barrier top does not supply it, for two
reasons: the electronic sector carries multiplicity three, so six states sit at
each Neumann level, and the half-line Neumann oscillator displaced from the origin
has low eigenvalues strictly below the full-line levels $1,3,5,\dots$. A
Golden-Thompson trace bound combined with Markov's inequality does supply it, in
the form
\begin{equation}
2k+1\le\frac{6}{\beta\Omega}\sqrt{\frac{1+x}{1-x}}\,
e^{\beta(\kappa B_{\mathrm{dimer}}+\Omega/2)},
\label{eq:count}
\end{equation}
and at $\beta=(\kappa B_{\mathrm{dimer}}+\Omega/2)^{-1}$ this gives $k+1\le11X$,
which is the count entering \eqref{eq:valleydoublets}. The prefactor is larger
than a counting argument tuned to the lowest pair would suggest, and it is stated
rather than optimized.

The all-temperature statement is the one that matters for a lift, because the
target strip is $T\ge A\Omega$ and a bound that degrades at $T\gtrsim\Omega$ would
be useless there.

\begin{theorem}[excursions at every temperature]\label{thm:valleyexcursion}
For every $\beta>0$, including $\beta\Omega\lesssim1$, and every
$\kappa\in(0,1)$,
\begin{equation}
\mathbb{P}_{\beta}\bigl(H_{r}>V_{\min}+\kappa B_{\mathrm{dimer}}\ \text{or}\
S=1\bigr)
\le\inf_{\gamma\in(0,1)}\frac{6}{\gamma}\sqrt{\frac{1+x}{1-x}}\,
e^{-\beta(1-\gamma)\kappa B_{\mathrm{dimer}}+\frac12\beta\gamma\Omega}
+3\,e^{-\beta(U+t^{2}/U)} .
\label{eq:valleyexcursion}
\end{equation}
In the same regime the pair-transfer overlap between any two sub-barrier
eigenstates decays as $O(e^{-c(x,\kappa)X})$ with no non-degeneracy assumption,
for cutoffs in $C^{2}(\mathbb{R};[0,1])$ the wall commutator obeys
$\|[H_{r},\chi_{R}]\psi_{k}\|_{L^{2}}\le C'_{\mathrm{com}}(x,\kappa)U\,
e^{-c(x,\kappa)X/2}$, the half exponent being the square root of the Agmon bound
on each local integral that enters it,
and crossing a single valley wall costs at least
$\beta(1-\gamma)\kappa B_{\mathrm{dimer}}$, that is $\Theta(\beta U)$, not
$\Theta(\beta\Omega)$.
\end{theorem}

The mechanism behind \eqref{eq:valleyexcursion} is an exact cancellation. The
Golden-Thompson upper bound on the relative partition function and the
constant-spinor lower bound both carry a factor $1/(\beta\Omega)$ from the
transverse oscillator, and the two cancel identically rather than being
controlled. This is precisely the $\beta\Omega$ bookkeeping that fails in the
phonon-number frame.

\subsection{The many-body barrier does not soften}

The step from the dimer to the lattice is a tensorization statement for the
Born-Oppenheimer potential
$V_{\mathrm{BO}}(v)=\sum_{x}v_{x}^{2}/(2U)-\operatorname{Tr}|h_{v}|-\mu|\Lambda|$.
It is the point at which one might fear that a collective distortion lowers the
barrier far below its dimer value, as it would for a soft mode. It does not.

\begin{theorem}[barrier tensorization]\label{thm:valleytensor}
Let $G=(\Lambda,E)$ have maximum degree $\le2d$ and let $\mu\in\mathbb{R}$. For
any $B\subset\Lambda$ and any distortion $v'$ with $v'=v$ on $B^{c}$ and
$|v'_{x}|=U$ on $B$,
\begin{equation}
V_{\mathrm{BO}}(v)-V_{\mathrm{BO}}(v')
\ge\sum_{x\in B}\Bigl[\frac{(U-|v_{x}|)^{2}}{2U}-4dt\Bigr].
\label{eq:valleytensor}
\end{equation}
\end{theorem}

The proof is a pinching inequality with the pinching family
$\{P_{B^{c}}\}\cup\{|x\rangle\langle x|\}_{x\in B}$, followed by a trace-norm
bound of $2t$ per edge through $\partial B$ and by the cancellation of the
unaffected electronic trace. Three consequences are worth separating. The per-site
barrier is at least $U/2-4dt$ for all $|v_{x}|\le U/2$, uniformly in the volume
and in $\mu$, so no collective path beats the dimer scale by more than the
coordination term. The statement is at the level of the potential, hence also at
the level of the quantum operator through the same inequality in quadratic form.
And the cold manifold never appears, so the route is not subject to the
counterexample of Section~\ref{sec:kernel}.

\subsection{Lang-Firsov against the Born-Oppenheimer exponent}

The distortion route predicts a splitting exponent set by the vibration barrier,
$c_{\mathrm{WKB}}(x,0)=2\sqrt{1-x^{2}}-2x^{2}\ln\frac{1+\sqrt{1-x^{2}}}{x}\in
[1.67856,\,2)$ for $x\in(0,1/4]$, while second-order Lang-Firsov polaron hopping
predicts
\begin{equation}
\Delta E_{\mathrm{LF}}^{(2)}\sim4\sqrt{\frac{\pi}{2U\Omega}}\,t^{2}\,
e^{-2U/\Omega},
\label{eq:lf2}
\end{equation}
that is $c_{\mathrm{LF}}=2$ with a $\sqrt{U/\Omega}$ prefactor and no barrier
correction, the standard small-polaron and bipolaron result
\cite{langfirsov,devreese2009}. The two exponents coincide only as $x\to0$. Exact oscillator-basis
diagonalization of the dimer agrees with \eqref{eq:lf2} in the antiadiabatic
regime $t^{2}/(U\Omega)\ll1$, where the ratio runs from $1.0007$ to $1.022$ over
$t\in[0.02,0.25]$ at $X=10$, and exceeds it once the adiabatic ratio
$t^{2}/(U\Omega)\gtrsim1$ sets in, reaching $7.658$ at $t/U=1/4$. The
Born-Oppenheimer exponent is the smaller of the two everywhere on $(0,1/4]$, which
is the correct qualitative ordering: the tunnelling particle is the distortion,
and the Lang-Firsov cloud is its large-$X$ limit.

\subsection{Where the route stops: locality in the distortion coordinate}

With Theorems~\ref{thm:valleyground}--\ref{thm:valleytensor} the dimer has a
splitting bound without the cold projection, and the lattice barrier cannot be
lowered collectively. The obstruction appears one step later and consists of three
interlocking facts.

First, the thermal energy is extensive. At temperature $T\ge A\Omega$ one has
$\langle H_{\Lambda}-E_{0}\rangle_{\beta}=\Theta(|\Lambda|A\Omega)$, which exceeds
$B_{\mathrm{dimer}}$ by a factor of order $|\Lambda|$; a global spectral
projection $\mathbf{1}_{\{H_{\Lambda}\le E_{0}+\kappa B\}}$ therefore carries
Gibbs weight $e^{-\Omega(|\Lambda|)}$ at strong coupling, and the dimer argument,
which uses exactly such a projection on $H_{r}$, has no lattice analogue.

Second, the local valley projectors do not commute with the boundary hopping. The
singlet weight $|1,1\rangle$ is $\Theta(t/U)$ at every distortion, and the local
phonon number is $\Theta(T/\Omega)=\Theta(A)$ in the target strip, so a local
projection onto the lower sheet leaks at first order in $t^{2}/U$ and not at
$e^{-cX}$. A finite witness makes this concrete: on the three-site $N=2$, $S=0$
electronic Hamiltonian at frozen distortion $v=(-U,+U,+U)$ with $U=20$, $t=2$, the
leakage of the neighbouring bond through the local ground-state projector is
\begin{equation}
\bigl\|(I-P_{01})V_{12}|\psi_{0}\rangle\bigr\|=0.1408944 ,
\label{eq:leak}
\end{equation}
close to $\sqrt{2}\,t^{2}/(2U)=0.14142$, and it is a power of $t$ rather than an
exponential in $X$. This is a fixed-distortion electronic witness, not a
three-site electron-phonon diagonalization; what it shows is that the projector
leakage inside a well is of the same order as the intra-well hopping that any
Schrieffer-Wolff step must resum, and there is no small parameter to resum in once
$\beta t^{2}/U\gg1$.

Third, inside a fixed valley the fermions remain dynamical. At $T\ge A\Omega$ and
$\Omega\le ct$ the neutral intra-valley amplitude obeys $\beta t^{2}/U\gg1$, so a
polymer expansion in $t$ is not small, and barrier crossings on adjacent bonds
couple through the fermion determinant at distance, with local electronic gaps
closing momentarily whenever $v_{x}(\tau)\approx\mu$. This is the same
determinant-coupling mechanism recorded in Section~\ref{sec:kernel} for the
charge-skeleton step, appearing now with the deformation coordinate in place of
the phonon number.

What survives is therefore a dimer-level lifting of the cold projection together
with a volume-uniform barrier bound, and what fails is the same
charged-skeleton influence bound as before, now asked of a distortion path rather
than of a phonon-number path. The route does not add an independent obstruction;
it relocates one, and it removes the cold projection from the part of the problem
where it was least harmful.

\section{The cutoff-free directed kernel, and why factorization fails}
\label{sec:kernel}

The failure of route (b) is informative, because the object that breaks
Gaussian domination, the influence action $\Phi_{\beta}$, can be integrated in
closed form along a charged transfer. The resulting kernel carries the desired
narrowing exponent, and it is the strongest positive object the programme has. It
is not enough, for a reason that is structural: an electron history is a sum, not
a product of independent transfers.

\subsection{The exact influence action and the physical pair insertion}

Work in the bare Lang-Firsov frame \eqref{eq:HLF} on the unrestricted phonon Fock
space. For a closed occupation history $q_{z}(u)=n_{z}(u)-1$ on the circle of
imaginary time, define
\begin{equation}
D_{\beta}(u)=\frac{\cosh\bigl(\Omega(\beta/2-u)\bigr)}{\sinh(\beta\Omega/2)},
\qquad K_{\beta}^{\mathrm{cov}}(u)=g^{2}D_{\beta}(u),
\label{eq:Db}
\end{equation}
so that $\int_{0}^{\beta}K_{\beta}^{\mathrm{cov}}(u)\,du=U$. Integrating the
driven oscillator gives
$Z_{\mathrm{ph}}\exp\bigl[\tfrac12\iint q_{z}K_{\beta}^{\mathrm{cov}}q_{z}\bigr]$,
and separating the Lang-Firsov atomic attraction gives the influence action
\begin{equation}
\Phi_{\beta}[q]=\frac U2\int_0^{\beta}q(u)^{2}du-\frac12\iint q K_{\beta}^{\mathrm{cov}}q
=\frac14\iint\bigl(q(u)-q(v)\bigr)^{2}K_{\beta}^{\mathrm{cov}}(u-v)\,du\,dv\ \ge0 .
\label{eq:Phi}
\end{equation}
An electron history with $h$ hops contributes its fermion sign, the factor
$t^{h}$, the ordered time integral, the odd-site atomic cost, and
$\exp\bigl[-\sum_{z}\Phi_{\beta}[q_{z}]\bigr]$. In Matsubara form, with
$\widehat q_{m}=\int q(u)e^{-i\omega_{m}u}du$ and jumps
$c_{j}=q(u_{j}^{+})-q(u_{j}^{-})$ summing to zero,
\begin{equation}
\Phi_{\beta}[q]=\frac{U}{2\beta}\sum_{m\ne0}\frac{\omega_{m}^{2}|\widehat q_{m}|^{2}}
{\Omega^{2}+\omega_{m}^{2}}
=\frac{U}{4\Omega}\sum_{j,k}c_{j}c_{k}D_{\beta}(u_{j}-u_{k}).
\label{eq:Phimodes}
\end{equation}
Since $D_{\beta}(u)$ has a minimum at $u=\beta/2$ and increases towards the
ends $u=0,\beta$, the difference $D_{\beta}(0)-D_{\beta}(u)$ is small near the
ends and largest at the antipode, which is the mechanism behind the coth in
\eqref{eq:dir2} below.

Two conventions matter and both are easy to get wrong. First, the physical pair
field in the transformed frame carries its own displacements,
\begin{equation}
U_{\LF}\Dop_{x}^{\dagger}\Dop_{y}U_{\LF}^{\dagger}
=\Dop_{x}^{\dagger}\Dop_{y}D_{x}(-2\alpha)D_{y}(2\alpha),
\label{eq:pairdisp}
\end{equation}
so that an equal-time pair correlation is not a trace over the undisplaced
Fock vacuum; dropping \eqref{eq:pairdisp} computes a different correlation. The
sign of both displacements may be reversed together without changing a Gaussian
trace, but not independently of the electron operator. Second, the atomic factor
and the influence action must both be kept: relative to the pair and empty
baselines, $\tfrac U2\int q^{2}=\tfrac U2\beta-\tfrac U2T_{\mathrm{odd}}$, and
neither the odd-site penalty nor $\Phi_{\beta}$ may be dropped.

\subsection{The free directed kernel}

Let $X_{zw}=D_{z}(\alpha)D_{w}(-\alpha)$ be the displacement product carried by a
coherent hopping process, and let $\langle\cdot\rangle_{\beta,\mathrm{ph}}$ be
the free oscillator trace. Normal ordering gives, for real $r,s$,
$\langle D(r;u)D(s;0)\rangle=\exp[-\tfrac12(r^{2}+s^{2})\coth(\beta\Omega/2)-rsD_{\beta}(u)]$,
and independence of the two sites then yields the exact cutoff-free identities
\begin{equation}
\begin{gathered}
F_{+}(u)=\bigl\langle X(u)X(0)\bigr\rangle,\qquad
F_{-}(u)=\bigl\langle X^{\dagger}(u)X(0)\bigr\rangle,\\
F_{\pm}(u)=\exp\Bigl[-\frac U\Omega\,
\frac{\cosh(\beta\Omega/2)\pm\cosh\bigl(\Omega(\beta/2-u)\bigr)}{\sinh(\beta\Omega/2)}\Bigr].
\end{gathered}
\label{eq:Fpm}
\end{equation}
Here $F_{+}$ is the directed weight and $F_{-}$ the neutral one. Since
$\cosh v\ge1$ and $(\cosh B+1)/\sinh B=\coth(B/2)$,
$F_{+}(u)\le e^{-X\coth(\beta\Omega/4)}$ for every $u\in[0,\beta]$, whence

\begin{theorem}[directed kernel bound]
\label{thm:dir}
For every $t$, $U$, $\Omega>0$ and $\beta>0$,
\begin{equation}
\Dir:=2t^{2}\int_{0}^{\beta}e^{-Uu}F_{+}(u)\,du
\le\frac{2t^{2}}{U}\bigl(1-e^{-\beta U}\bigr)\,e^{-X\coth(\beta\Omega/4)} .
\label{eq:dir2}
\end{equation}
\end{theorem}

\begin{proof}
Insert $F_{+}(u)\le\exp[-X\coth(\beta\Omega/4)]$ into the integral and evaluate
$\int_{0}^{\beta}e^{-Uu}du$.
\end{proof}

The exponent in \eqref{eq:dir2} is the right one: at $\beta\Omega\gg1$ it is
$X[1+O(e^{-\beta\Omega})]$, at $\beta\Omega\ll1$ it is $4X/(\beta\Omega)$, which
is $4U\tau/\Omega^{2}$, the classical narrowing, and it interpolates between them
monotonically. It is also the same function that appears in the released
cold-projected exponent.

What Theorem~\ref{thm:dir} is \emph{not} is a bound on an interacting Gibbs
activity. It is a statement about a free oscillator trace with one inserted
charge pair, and it says nothing about how much the surrounding electron history
costs or gains. The remainder of this section is about exactly that.

\subsection{Conditional factorization is false}

The natural way to use Theorem~\ref{thm:dir} is to bound the weight of an
electron history by the product of the weights of its charged transfers, in other
words to assume that the Gaussian factor factorizes over disjoint time segments.
Precisely, for a chronological sequence of real bond charges $k_{j}$ at times
$u_{j}$, whether or not $\sum k_{j}=0$, the free oscillator trace over the two
endpoints is
\begin{equation}
W(k,u)=\exp\Bigl[-\alpha^{2}\sum_{j,k}k_{j}k_{k}D_{\beta}(u_{j}-u_{k})\Bigr],
\label{eq:W}
\end{equation}
and for a division of the charges into two groups $A$ and $B$ the ratio
$W(A\cup B)/[W(A)W(B)]$ has logarithm
$-2\alpha^{2}\sum_{j\in A,k\in B}k_{j}k_{k}D_{\beta}(u_{j}-u_{k})$, which has no
fixed sign. Positivity of the full quadratic form does not imply positivity of its
cross term, and hence not the existence of independent activities.

\begin{proposition}[explicit neutral enhancement]
\label{prop:witness}
Let $U=40$, $\Omega=1$, $\beta=1$ and $\epsilon=1/160$. Take the four bond charges
$(+1,+1,-1,+1)$ at the ordered times $(\epsilon,2\epsilon,3\epsilon,1/2)$,
starting from a doublon and ending in a pair that has been transferred. Every
intermediate electronic occupation is legal, namely $(0,2),(1,1),(2,0),(1,1)$,
and the last two hops are a completed neutral excursion around the post-transfer
state. With $W_{4}=W\bigl((1,1,-1,1),u\bigr)$ and $W_{N}=F_{-}(1/2-3\epsilon)$,
\begin{equation}
\log\frac{W_{4}}{W_{N}}+40\coth(1/4)\in[9.295999605219290,\ 9.295999605219293].
\label{eq:witness}
\end{equation}
In particular $\log(W_{4}/W_{N})>-40\coth(1/4)$, so $W_{4}/W_{N}$ exceeds the
largest value that any free directed block $F_{+}(u)$ can take at these
parameters, not merely $F_{+}(\epsilon)$.
\end{proposition}

The interval in \eqref{eq:witness} is an enclosure computed in fixed-precision
interval arithmetic from the closed form \eqref{eq:Db}; the two endpoints differ
by $3\times10^{-15}$. The atomic odd-site durations split additively between the
two groups, so multiplying by the atomic factors does not repair the
factorization. This is a counterexample to a conditional local Gaussian bound,
not to an integrated and suitably reorganized estimate.

The witness is not an accident of $X=40$. It extends to an analytic family in the
whole strong-coupling regime.

\begin{theorem}[a persistent failure at $\beta\Omega=1$]
\label{thm:family}
Fix $\beta\Omega=1$, let $X\ge20$ and $h=1/(4X)$. Let $D(v)=\cosh(1/2-v)/\sinh(1/2)$
and use the dimensionless times $(h,2h,3h,1/2)$ with charges $(+1,+1,-1,+1)$, as
in Proposition~\ref{prop:witness}. Then
\begin{equation}
\log\frac{W_{4}}{W_{N}}+X\coth(1/4)\ \ge\ \frac{X}{5}-\frac{13}{4}\ \ge\ \frac34>0 .
\label{eq:family}
\end{equation}
\end{theorem}

\begin{proof}
On $[0,1]$, $|D'(v)|\le1$. At $h=0$ the three coincident unit charges combine to
$+1$, so that $\log W_{4}$ and $\log W_{N}$ are evaluated with two unit charges
separated by $\beta/2$. Using $(\cosh(1/2)+1)/\sinh(1/2)=\coth(1/4)$ and
$(\cosh(1/2)-1)/\sinh(1/2)=\tanh(1/4)$,
\[
\log W_{4}=-2\alpha^{2}\coth(1/4)=-X\coth(1/4), \qquad
\log W_{N}=-2\alpha^{2}\tanh(1/4)=-X\tanh(1/4),
\]
while the free directed maximum is $-X\coth(1/4)$ by \eqref{eq:Fpm}. The margin at
$h=0$ is therefore exactly $X\tanh(1/4)$. Moving the first three charges to
distinct legal times shifts $\log W_{4}$ by at most $10Xh$, the six unordered
separations being $h,2h,h,h,2h,3h$ with sum $10h$, and shifts $\log W_{N}$ by at
most $3Xh$, because each exponent is $X$ times a signed sum of nonnegative
$D$ values with $|D'|\le1$. Hence the legal margin is at least
$X\tanh(1/4)-13/4\ge X/5-13/4\ge3/4$ for $X\ge20$, where $\tanh(1/4)\ge1/5$
follows from $e^{1/2}\ge1+1/2$.
\end{proof}

Choosing $\Omega=\kappa t^{2}/U$ with fixed $\kappa>0$ and $t/U$ small places this
family at $\tau=\Omega$ inside the strong-coupling temperature window of
\cite{largeu,strong}, whenever $\kappa$ is below the constant there. Its long
neutral odd excursion is atomically costly, which is why it is a counterexample to
the conditional local bound and not to a normalized integrated ceiling.

\subsection{A closed pair-return loop costs only \texorpdfstring{$O(1)$}{O(1)}}

\begin{proposition}[short closed loops]
\label{prop:loop}
Take four hops on a bond with charges $(+1,+1,-1,-1)$ at $(0,\epsilon,3\epsilon,4\epsilon)$,
$\epsilon=1/(4U)$, and $\beta>4\epsilon$. The occupations are
$(0,2),(1,1),(2,0),(1,1),(0,2)$, a closed and legal electron history. Then
\begin{equation}
\Phi_{\beta}[q]\le\tfrac52,\qquad \tfrac U2T_{\mathrm{odd}}=\tfrac12,
\qquad\text{so the total action is at most }3,
\label{eq:loop}
\end{equation}
independently of $U/\Omega$ as $\Omega/U\to0$ with $\beta\Omega=1$.
\end{proposition}

\begin{proof}
At each of the two sites, drop the nonnegative covariance form after subtracting
the static baseline: $\Phi[q]=\Phi[q-q_{0}]\le\tfrac U2\int(q-q_{0})^{2}$, where
$q_{0}$ is the time average. With the durations $T=3\epsilon$ at each site this
gives $\Phi\le U(4T-2\epsilon)=\tfrac52$, and the two odd-site atomic costs
$2U\epsilon$ combine to $\tfrac12$.
\end{proof}

Both directed free weights in this history are exponentially small in $X$, while
the closed loop's weight is at least $e^{-3}$. A short virtual loop does not by
itself force pair ordering, but it does mean that any estimate whose activities
are the charged weights of Section~\ref{sec:kernel} undercounts the true weight of
the history that carries a pair by an $O(1)$ factor at every return.

\subsection{The endpoint insertion has the same short-time behaviour}

A site with $n=2$ for a duration $T$ and $n=0$ otherwise has jumps $+2,-2$ and
exact action $\Phi_{\mathrm{site}}(T)=2X[D_{\beta}(0)-D_{\beta}(T)]$, while a
pair of opposite sites has $\Phi_{\mathrm{pair}}(T)=4X[D_{\beta}(0)-D_{\beta}(T)]$.
At $T=\beta/2$ the site action is $2X\tanh(\beta\Omega/4)$, not $4X\tanh(\beta\Omega/4)$,
and for fixed $\beta\Omega>0$ and $T=1/U$ the site action tends to $2$, not $4$.
At $\beta\Omega\ll1$ the antipodal site action is $\beta U/2+O(\beta U(\beta\Omega)^{2})$,
not $\beta U$. The coth of the open kernel therefore cannot be interchanged with
the tanh of a closed loop, and any estimate that does so misstates the action by a
constant factor at short times, which in a cluster expansion is the difference
between convergence and divergence.

\subsection{A facing-plateau lemma, and what it does not do}

One local positive statement does survive on the interacting side. It applies to
a restricted geometry, which is part of its hypotheses and not a consequence of
the model.

\begin{theorem}[facing plateaus]
\label{thm:plateau}
Assume facing sign-restricted intervals of length at least $L_{0}$, separated by
an odd-only gap of length $\delta$, with the intervals and gap fitting on the
circle without overlap: on the left interval $q\in\{-1,0\}$, on the right
$q\in\{0,1\}$, and $q=0$ on the gap. Let $W=\min(L_{0},1/\Omega)$ and let
$T_{\mathrm{odd}}=|\{q=0\}|$. Then
\begin{equation}
\frac U2T_{\mathrm{odd}}+\Phi[q]\ \ge\ \frac{3}{22}\,UW .
\label{eq:plateau}
\end{equation}
\end{theorem}

\begin{proof}
If $T_{\mathrm{odd}}\ge W$ the atomic term alone suffices. Otherwise
$\delta\le T_{\mathrm{odd}}<W$. Take a trapezoid $g$ equal to one on the gap, with
ramps of width $W$ on each facing interval and zero outside. Cauchy-Schwarz over
the nonzero Matsubara modes, using \eqref{eq:Phimodes}, gives
$\Phi[q]\ge\tfrac U2\,|\!\int qg'|^{2}/\int(\Omega^{2}g^{2}+(g')^{2})$. The
numerator's square root is at least $2-T_{\mathrm{odd}}/W>1$, and the denominator is
$2/W+\Omega^{2}(2W/3+\delta)<11/(3W)$. Hence $\Phi[q]\ge3UW/22$.
\end{proof}

Two limitations decide the fate of the plateau route. The facing geometry is not
forced: two arbitrary disjoint intervals do not produce it, and a charged
transfer is not guaranteed to possess such intervals at all. Even where it does,
several transfers may share the same cost, and the statement has not been shown
additive along a spatial skeleton. A local gain of $3UW/22$ per plateau is worth
nothing in a long-distance estimate unless the costs can be allocated to distinct
transfers with bounded overlap and unless the dwell time that produces the
plateau is itself controlled. Neither is available.

\subsection{What the kernel route still needs}

The failures of this section narrow the target. A candidate decomposition must
distinguish short virtual transfers, whose spatial activity is algebraically
small, from long-lived pair transport, and for the latter it must prove a
\emph{normalized charged-skeleton} statement rather than multiply free $F_{+}$
traces. Section~\ref{sec:openproblem} states that target precisely.

\section{The hot sector in the full Gibbs state}
\label{sec:vtail}

\subsection{What is measured}

Sections~\ref{sec:routes} and~\ref{sec:kernel} describe estimates that fail. This
section describes one that succeeds, and it is the only volume-uniform positive
statement in the note. Work in the bare Lang-Firsov frame of
Section~\ref{sec:model}, with the frozen cutoff
\begin{equation}
M=\Bigl\lfloor\frac{X}{64}\Bigr\rfloor,\qquad
Q_{z}=\mathbf{1}_{\{N_{z}\le M\}},\qquad F_{z}=\mathbf{1}-Q_{z},
\qquad p=\operatorname{Tr}\bigl(\varrho F_{z}\bigr),
\label{eq:Fz}
\end{equation}
where $\varrho=Z^{-1}e^{-\beta H}$ is the full Gibbs state of the transformed
Hamiltonian. Because the Lang-Firsov transformation is a unitary conjugation, $p$
is the Gibbs expectation of an honest observable; nothing about it depends on
choosing a frame. What $p$ is not is the Gibbs weight of the complement of the
cold projection used in \cite{largeu}: that object does not appear here. Nor is
$p$ computed in a truncated model. The Gibbs state is the exact one, on the
unbounded Fock space, and the phonon number is never cut in the destination
space; the cutoff enters only as the target of one local bounded operator, whose
matrix elements are estimated by Cauchy contours.

Two logically distinct uses of an upper bound on $p$ occur in the literature on
this problem. In a \emph{witness} route one hopes to compare the Gibbs state with
its cold-projected image, which requires the complement to carry exponentially
small weight. In an \emph{activity} route one hopes that the hot sector behaves
like a rare defect in a polymer expansion, which requires its activity to be
exponentially small in $X$. The theorem below shows that neither hope can be
granted at the level of the total weight: the complement is always at least a
fixed power of $t/U$.

\subsection{Two theorems, and the local commutator proof}

\begin{theorem}[dimer]\label{thm:vtailD}
Fix $0<a<b<\infty$ and put
\[
c_{0}=\min\Bigl\{\frac1{64},\frac{a}{64+16\ln2}\Bigr\},\qquad
c_{1}=\min\{1,1/b\},\qquad c_{2}=\frac1{4096}.
\]
On the two-site particle-hole-symmetric spinful Holstein dimer, if
$a\le\beta\Omega\le b$, $0<\Omega<c_{0}U$ and $0<t<c_{1}\Omega^{3}/U^{2}$, then
$\operatorname{Tr}(\varrho F_{z})\ge c_{2}(t/U)^{2}$ at either site.
\end{theorem}

\begin{theorem}[volume-uniform]\label{thm:vtailT}
Let $\Lambda$ be a finite nearest-neighbour torus $(\mathbb{Z}/L\mathbb{Z})^{d}$
with $d\ge1$, $L\ge3$, and $q=2d$; count each undirected edge once. With
\[
c_{0}=\min\Bigl\{\frac1{64},\frac{a}{64+16\ln(q+1)}\Bigr\},\qquad
c_{1}=\min\Bigl\{\frac1{2^{20}q^{3}},\frac1{bq}\Bigr\},\qquad
c_{2}=\frac1{4096},
\]
the conclusion of Theorem~\ref{thm:vtailD} holds at every site, with no constant
depending on $L$. The same conclusion holds on every finite vertex-transitive
simple graph of degree $q$.
\end{theorem}

The proof is a local commutator computation, and its structure is worth recording
because it shows exactly where the phonon tail is created.

\emph{Step 1: cold displacement decay.} For the one-mode cold projector
$Q=\mathbf{1}_{\{N\le M\}}$ and $\gamma=\alpha$ or $2\alpha$, two Cauchy contours
of radius $\sqrt M$ applied to the Bargmann kernel of $D(\gamma)$ give
\begin{equation}
\bigl\|QD(\gamma)Q\bigr\|\le\delta(X),\qquad
\delta(X)=\Bigl(\frac{X}{64}+1\Bigr)e^{-c_{D}X},\qquad
c_{D}=\frac14-\frac1{4\sqrt{2}}-\frac1{64}>0 ,
\label{eq:delta}
\end{equation}
for $X\ge64$. Hence on two cold sites the two opposite displacement images
overlap in norm by at most $\delta^{2}$.

\emph{Step 2: a positive virtual block.} For one edge $e=\{z,w\}$, let $P$ project
onto the cold pair and empty arrangements of the edge, $R=P_{2}F_{z}$ with $P_{2}$
the two-odd-site sector, $L=U+\Omega(N_{z}+N_{w})$ and
$\mathcal T=RL^{-1}V_{e}P$, and set $A=2PV_{e}RL^{-1}V_{e}P$. Then
\begin{equation}
A\ \ge\ \frac23\,\frac{t^{2}}{U}\,P .
\label{eq:Apos}
\end{equation}
The rank of $P$ is finite while $R$ has infinite rank: $A$ is an operator of
finite rank whose range contains the phonon tail, and \eqref{eq:Apos} is a
strictly positive virtual second-order amplitude. The proof uses only the
two-dimensional charge component sum, \eqref{eq:delta}, and the
Cauchy--Schwarz inequality in the $L^{1/2}$ and $L^{-1/2}$ norms.

\emph{Step 3: a local Ehrenfest identity.} With $K=\mathcal T-\mathcal T^{\dagger}$
one has, exactly on the edge,
\begin{equation}
[K,H_{0}]=-B-B^{\dagger}+\mathcal T E_{c}+E_{c}\mathcal T^{\dagger},
\qquad \|E_{c}\|\le U/32,
\label{eq:ehrenfest}
\end{equation}
and every term on the right except the first two is bounded by $t\sqrt p$.
Using trace Cauchy--Schwarz for operators with one endpoint in a hot or odd
projector, and the trace identity $\operatorname{Tr}(\varrho[K,H])=0$, one
obtains
\begin{equation}
\frac23\langle P\rangle\le(5+20x)\frac{\sqrt p}{x}+32(q-1)\sqrt r ,
\label{eq:main}
\end{equation}
where $x=t/U$ and $r=\operatorname{Tr}(\varrho\,o_{z})$ is the probability that a
site is odd. The trace identity needs care on the unbounded Fock space: $K$ is not
bounded, but $\mathcal T$ satisfies $H_{0,e}\mathcal T=RV_{e}P$ and
$\mathcal T H_{0,e}=\mathcal T E_{c}$ locally and preserves
$\operatorname{Dom}(H_{0})$, in the sense that
$\mathcal T\operatorname{Dom}(H_{0})\subseteq\operatorname{Dom}(H_{0})$. On a
torus with sites outside the edge, $\mathcal T$ does not map the whole Hilbert
space into $\operatorname{Dom}(H_{0})$, because the exterior energy stays
unbounded; the exterior part of $H_{0}$ commutes with $\mathcal T$ and cancels in
the commutator on that domain. Hence $K$ preserves $\operatorname{Dom}(H_{0})$
and $[K,H_{0}]$ extends from it to the bounded operator in \eqref{eq:ehrenfest}.
Expanding in the discrete eigenbasis of the bounded-below $H$ then justifies the
identity term by term.

\emph{Step 4: two local inputs.} Removing the $q$ incident edges at a site changes
the Hamiltonian in norm by at most $2qt$, and concavity of the local free energy
in a field coupling to $o_{z}$ gives
\begin{equation}
r\le\frac{4}{\beta U}e^{-\beta U/4}+16qx .
\label{eq:rbound}
\end{equation}
Second, the on-site unitary that interchanges $|0\rangle$ and $|2\rangle$ changes
only the edges incident to one endpoint, so by quantum Pinsker
$|\langle E\rangle-\langle E_{\mathrm{eq}}\rangle|\le\sqrt{2q\beta t}$, and since
$E+E_{\mathrm{eq}}=\mathbf{1}-o_{z}-o_{w}$,
\begin{equation}
\langle P\rangle\ge\tfrac12-r-\sqrt{q\beta t/2}-2p .
\label{eq:Pbound}
\end{equation}
Inserting \eqref{eq:rbound} and \eqref{eq:Pbound} into \eqref{eq:main} and using
$X>64+16\ln(q+1)$ closes the constants. The two logically separate cases
$p\ge1/16$ and $p<1/16$ both end at $p>x^{2}/4096$.

Three features of this proof are the point. It uses no positivity of fermionic
worldline weights, so it applies to the spinful model with signs. It pays
explicitly for the absence of a local particle-hole symmetry at nonzero hopping;
equality of paired and equal-even probabilities would be false. And every
estimate is local, so the constants do not move with the volume: only the edge and
its $2(q-1)$ neighbours enter. What the theorem does \emph{not} do is bound the
effect of the hot sector on pair propagation. It measures a weight, not an
influence.

\subsection{An independent second-order route on the dimer}

A second, logically independent derivation reaches the same conclusion on the
dimer by a Duhamel expansion, with different constants and a much narrower hopping
window. Write $V=-tW$ with $\|W\|=2$ and consider the numerator
$N(t)=\operatorname{Tr}(F_{z}e^{-\beta(H_{0}-tW)})$. Hopping has no diagonal
electronic matrix elements, and since $F_{z}$ and $H_{0}$ are simultaneously
diagonal every second-order simplex summand is nonnegative. Keeping only the
channels that start in a hot two-single-occupancy singlet and visit a pair vacuum
gives, after the exact simplex integral
$\int_{0}^{\beta}(\beta-s)e^{sE}ds=(1-(1+\beta E)e^{-\beta E})/E^{2}$ and a
Poisson concentration estimate on the Franck-Condon sum,
\begin{equation}
\frac{t^{2}N_{2}}{Z_{0}}\ge\frac{h}{256}\,x^{2},
\qquad h=(1-e^{-a})^{2}.
\end{equation}
The all-orders remainder obeys
$|R_{\ge3}|\le\frac43Z_{0}(\beta t)^{3}e^{2\beta t}$, by noncommutative
Schatten--H\"older on each Duhamel simplex, so that with
\begin{equation}
c_{0}^{\mathrm{Du}}=\min\Bigl\{\frac1{64},\frac{a}{16}\Bigr\},\qquad
c_{1}^{\mathrm{Du}}=\min\Bigl\{\frac1b,\frac{h}{2^{15}b^{3}}\Bigr\},\qquad
c_{2}^{\mathrm{Du}}=\frac{h}{1024}
\label{eq:duconstants}
\end{equation}
the same hypotheses give $N(t)/Z(t)\ge c_{2}^{\mathrm{Du}}x^{2}$ whenever
$\beta t\le b\,c_{1}^{\mathrm{Du}}/X^{2}$. These constants are sufficient, not
optimized. The window has the same $t=O(\Omega^{3}/U^{2})$ scaling as
Theorem~\ref{thm:vtailD} with a much smaller prefactor $c_{1}^{\mathrm{Du}}$, and a
volume-extensive remainder of this type would
\emph{not} give Theorem~\ref{thm:vtailT}: the local argument is what makes the
extension volume-uniform.

\subsection{The channel coefficient, and what the first rotation does not do}

It is worth separating the part of the virtual cloud that comes from the
Lang-Firsov denominator from the part that is Franck-Condon in origin. Let
$p_{X}(n)=e^{-X/2}(X/2)^{n}/n!$ and define the dimensionless channel weights
\begin{align}
w_{\mathrm{LF}}(X)&=\sum_{n_{z}>M,\,n_{w}\ge0}
\frac{p_{X}(n_{z})p_{X}(n_{w})}{\bigl(1+(n_{z}+n_{w})/X\bigr)^{2}},
\label{eq:wLF}\\
w_{S_{1}}(X)&=\sum_{n_{z}>M,\,n_{w}\ge0}p_{X}(n_{z})p_{X}(n_{w})
\Bigl(\frac{(n_{z}+n_{w})/X}{1+(n_{z}+n_{w})/X}\Bigr)^{2}.
\label{eq:wS1}
\end{align}
The second line is the same channel after the first Schrieffer-Wolff rotation of
\cite{largeu}, whose generator removes the electronic gap $U$ from the
denominator: the amplitude $1/(1+s/X)$ becomes $1/(1+s/X)-1=-s/X/(1+s/X)$. At the
Franck-Condon peak $s=n_{z}+n_{w}\sim X$ this changes $+1/2$ into $-1/2$, and the
squared weight is unchanged. Poisson concentration gives
\begin{equation}
\lim_{X\to\infty}w_{\mathrm{LF}}(X)=\lim_{X\to\infty}w_{S_{1}}(X)=\frac14,
\qquad
w_{\mathrm{LF}}=\frac14+\frac{3}{16X}+O(X^{-2}),\quad
w_{S_{1}}=\frac14-\frac{1}{16X}+O(X^{-2}).
\label{eq:wlimit}
\end{equation}
So the tail is not an artifact of the bare-frame denominator. The leading
Franck-Condon channel keeps a coefficient of order one in the strong-coupling
limit, and the rotation specified in \cite{largeu} removes the electronic energy
cost but not the phonon energy of order $U$. An energy-resolving generator can
cancel each nonresonant vacuum channel at first order, but it cannot be extended
across the exact resonances of the full Fock space. \emph{A nonperturbative lower
bound for the rotated full Gibbs state is open}; \eqref{eq:wlimit} is a statement
about a specified perturbative channel, not about that state.

\subsection{How to read a positive lower bound}

Three readings, in increasing order of what is actually supported.

First, \eqref{eq:wlimit} and Theorems~\ref{thm:vtailD} and~\ref{thm:vtailT} do not
falsify the lifted ceiling, and nothing in this note does. A ceiling is an upper
bound on a temperature, and it is entirely compatible with rapid pair
decorrelation, with a hot sector of algebraic weight, and with the failure of
every factorization attempted in Section~\ref{sec:kernel}. The cold-projected
theorem of \cite{largeu} is not in question.

Second, the results do constrain the shape of any proof. An argument that expels
the hot sector by showing it is exponentially rare cannot exist: the weight is
bounded below by $c_{2}(t/U)^{2}$ at every temperature in the stated window, with
constants that do not depend on the volume. Any comparison between the Gibbs
state and its cold-projected image must therefore tolerate an algebraically large
displaced component, which is why the comparison is naturally a statement about
influence and not about weight.

Third, the lower bounds are compatible with the upper bounds one gets from the
influence action. The exact action
$\Phi_{\beta}[q]=\frac{U}{4\Omega}\sum_{j,k}c_{j}c_{k}D_{\beta}(u_{j}-u_{k})$ is
nonnegative on closed charge histories, so a large total weight can coexist with
strong suppression of every individual long-lived decoration. That coexistence is
precisely the reason a bound on the total weight is not enough; the missing object
is described next.

\section{Charge order: the same wall seen from the other side}
\label{sec:cdw}

\subsection{The strong-coupling pseudospin theory, as known results}

Half-filled models of this type have a second strong-coupling limit in which the
electron pair on a site plays the role of a hard-core boson. It is worth recording
what is known in that limit, because the same wall that obstructs the pairing
ceiling obstructs the ordering theorem there, and because the pseudospin picture
is sometimes offered as an alternative route to a ceiling.

On a periodic hypercubic torus with $N=L^{d}$, $d\ge2$, $z=2d$, consider
\begin{equation}
H=-t\sum_{\langle x,y\rangle,\sigma}\bigl(c^{\dagger}_{x\sigma}c_{y\sigma}
+\mathrm{h.c.}\bigr)+\sum_{x}\Bigl[\frac{p_{x}^{2}}{2m}+\frac K2q_{x}^{2}
+gq_{x}(n_{x}-1)\Bigr],
\label{eq:cdwH}
\end{equation}
with $U=g^{2}/K$, $\omega=\hbar\Omega=\hbar\sqrt{K/m}$, $X=U/\omega$ and
$\alpha^{2}=X/2$. The shift $n_{x}-1$ makes the model particle-hole symmetric at
half filling. After the Lang-Firsov transformation and subtraction of the
oscillator zero point and the atomic ground energy,
$H_{\mathrm{LF}}=H_{0}+Q$ with $h_{x}=\omega b_{x}^{\dagger}b_{x}
+\frac U2\mathbf{1}_{\{n_{x}=1\}}$ and bond hopping dressed by opposite
displacements. The dressed hopping has norm exactly $2t$.

Second-order degenerate perturbation theory in the cold pair-hole subspace gives
an easy-axis XXZ pseudospin model with
\begin{equation}
f_{z}(X)=e^{-X}\,{}_{1}F_{1}(X;X+1;X),\qquad
f_{xy}(X)=e^{-X}\,{}_{1}F_{1}(X;X+1;-X)
=e^{-2X}{}_{1}F_{1}(1;X+1;X),
\label{eq:fz}
\end{equation}
$J_{a}=4t^{2}f_{a}/U$ and
\begin{equation}
H_{2}=\sum_{b}\Bigl[J_{z}\bigl(\tau_{x}^{z}\tau_{y}^{z}-\tfrac14\bigr)
-\frac{J_{xy}}{2}\bigl(\tau_{x}^{+}\tau_{y}^{-}
+\tau_{x}^{-}\tau_{y}^{+}\bigr)\Bigr],
\label{eq:H2}
\end{equation}
with no additional constant outside the bond sum. In the integral form used
below,
\begin{equation}
f_{z,\;xy}(X)=Xe^{-X}\!\int_{0}^{1}\!u^{X-1}e^{\pm Xu}\,du,
\qquad
f_{z}(X)-f_{xy}(X)=2Xe^{-X}\!\int_{0}^{1}\!u^{X-1}\sinh(Xu)\,du>0 .
\label{eq:fdiff}
\end{equation}

Every statement in \eqref{eq:fz}--\eqref{eq:fdiff} is a known result, and this
note claims none of it. The bipolaron or pseudospin effective Hamiltonian of the
strong-coupling electron-phonon problem, with distinct pair-transfer and
pair-hole diagonal weights and hence easy-axis anisotropy, was derived by
Alexandrov and Ranninger \cite{ar81} and by Hirsch and Fradkin \cite{hf83}, who
mapped the strong-coupling molecular-crystal Holstein model onto this
spin-$\frac12$ XXZ chain, isotropic as $\omega\to\infty$ and Ising as
$\omega\to0$. The review of
Micnas, Ranninger and Robaszkiewicz \cite{mrr90} describes the same mapping and
the resulting competition between charge and pair coherence. The exact integral
representation in \eqref{eq:fdiff} is published: Han, Kivelson and Yao
\cite{hky20}, in the supplement of their paper, give
$F(\pm X,Y)=Ye^{-X}\int_{0}^{1}z^{Y-1}e^{\pm Xz}dz$, which at $Y=X$ is
\eqref{eq:fdiff}. Freericks \cite{freericks93} rederived the second-order
coefficients as functions of $X$ and extended the expansion. The elementary
consequences are recorded here only for completeness: $J_{z}\to 2t^{2}/U$ and
$J_{xy}=O(\sqrt{X}e^{-2X})$ as $X\to\infty$, both tend to $4t^{2}/U$ as
$X\to0$, and the anisotropy is positive at every finite frequency.

\subsection{What this note adds: a full-Fock transform and two trace lemmas}

The proofs in this area usually proceed through an effective Hamiltonian. What is
added here is that the first step can be done on the unrestricted Fock space, with
volume-uniform constants and no phonon cutoff, and that two of the trace inputs
needed later can be stated and proved in the same setting.

\emph{A bounded Lie-Schwinger generator.} For a bond $b=\{x,y\}$ let
$P_{b}=P_{x}P_{y}$ with $P_{x}=p_{x,-}+p_{x,+}$ the cold pair-hole projectors,
and set
\begin{equation}
A_{b}=(H_{0,b}|_{P_{b}^{\perp}})^{-1}Q_{b}P_{b},\qquad
S_{b}=A_{b}-A_{b}^{\dagger},\qquad
O_{b}=P_{b}Q_{b}+Q_{b}P_{b},\qquad
R_{b}=(\mathbf{1}-P_{b})Q_{b}(\mathbf{1}-P_{b}).
\end{equation}
Although $(H_{0,b}|_{P_{b}^{\perp}})^{-1}$ has norm $1/D$ with
$D=\min(U/2,\omega)$, its action on $Q_{b}P_{b}$ lands in the two-single-electron
sector where $H_{0,b}\ge U$, so that
$\|A_{b}\|=\|S_{b}\|\le2t/U$, the rank of $A_{b}$ on the two-site bond space is at
most two, that of $S_{b}$ at most four, and
$\|O_{b}\|,\|R_{b}\|\le2t$. Each $S_{b}$ preserves $\operatorname{Dom}(H_{0})$,
the finite-volume sum $S=\sum_{b}S_{b}$ is bounded and anti-self-adjoint, and the
transformation
\begin{equation}
e^{S}H_{\mathrm{LF}}e^{-S}=H_{0}+R+
\sum_{n\ge1}\Bigl[\frac1{n!}\mathrm{ad}_{S}^{n}R
+\frac{n}{(n+1)!}\mathrm{ad}_{S}^{n}O\Bigr]
\label{eq:LS}
\end{equation}
holds exactly on $\operatorname{Dom}(H_{0})$ in finite volume, with
$O=\sum_{b}O_{b}$ and $R=\sum_{b}R_{b}$. The closedness argument is two lines:
for $\psi\in\operatorname{Dom}(H_{0})$,
$\|H_{0}S^{k}\psi\|\le\|S\|^{k}\|H_{0}\psi\|+k\|S\|^{k-1}\|O\|\|\psi\|$, so the
exponential series converges in the graph norm.

\emph{Connected-word bounds.} Nonoverlapping even fermionic operators commute, so
a nested commutator with a root bond has a connected support of at most $k+2$
sites at step $k$, with at most $z(k+2)$ choices of the next bond. With
$r=4zt/U<1$ the terms of \eqref{eq:LS} from three hopping events on obey
\begin{equation}
\frac{\|E_{\ge3}\|}{N}\le zt\,\frac{r^{2}(5-3r)}{(1-r)^{2}},\qquad
\|E_{\ge3}\|_{2}\le4zt\,\frac{q^{2}(20-25q+9q^{2})}{(1-q)^{3}},
\label{eq:cw1}
\end{equation}
for $q=2r<1$, where the exponentially weighted norm is
$\|\Phi\|_{2}=\sup_{x}\sum_{B\ni x}2^{|B|}\|\Phi_{B}\|$ for the interaction
decomposition induced by connected words. On a bipartite lattice the operator
$G=(-1)^{N_{A}}$ anticommutes with $Q_{b}$ and $S_{b}$ and commutes with $P$, so
$P[S,R]P=0$ and the cold compression of the cubic term vanishes; the compressed
remainder begins at four hopping events, with
\begin{equation}
\frac{\|PE_{\ge3}P\|}{N}\le zt\,\frac{r^{3}(7-3r^{2})}{(1-r^{2})^{2}},\qquad
\|PE_{\ge3}P\|_{2}\le4zt\,\frac{q^{3}(35-28q^{2}+9q^{4})}{(1-q^{2})^{3}} .
\label{eq:cw2}
\end{equation}

\emph{An explicit local window.} In the window
\begin{equation}
\frac U8\le\omega\le\frac U4,\qquad \frac tU\le\frac1{65536\,d^{2}},
\label{eq:window}
\end{equation}
one has $X\in[4,8]$ and $r\le1/16384$. Jensen's inequality for the Poisson mean
gives $f_{z}(X)>1/2$, while $0<f_{xy}(X)<e^{-X}\le e^{-4}<3/103$, using the exact
witness $e^{4}>103/3$. Hence
\begin{equation}
\frac{J_{xy}}{J_{z}}<\frac{6}{103}<\frac1{16},\qquad
\Delta:=J_{z}-J_{xy}>\frac{194}{103}\frac{t^{2}}{U},\qquad
\|PE_{\ge3}P\|_{2}<\frac{\Delta}{2048},
\label{eq:window2}
\end{equation}
the last comparison being exactly $17\cdot103<194\cdot16$. All constants in
\eqref{eq:window2} were verified in exact rational arithmetic, and the
coefficient values in \eqref{eq:fz} were enclosed in rational intervals narrower
than $10^{-40}$ at $X=\frac1{16},\frac14,1,2,4,8,16$ by two independent
implementations using different representations.

\emph{A fixed-history trace bound.} Let $D=\min(U/2,\omega)$. On a single site
\begin{equation}
\operatorname{Tr}_{\mathcal H_{x}}\bigl(Q_{x}^{h}e^{-vh_{x}}\bigr)
=\frac{2\bigl(e^{-vU/2}+e^{-v\omega}\bigr)}{1-e^{-v\omega}}
\le\frac{4e^{-vD}}{1-e^{-vD}},
\label{eq:htrace}
\end{equation}
and for $vD\ge4$ this is at most $e^{-vD/2}$, because
$4e^{-2}+e^{-4}<21/25<1$ is an exact witness. Consequently, for
$\beta D\ge4$, for any split $0<\ell_{1},\dots,\ell_{k}$ with
$\sum_{i}\ell_{i}=\beta$ and for bounded operators $B_{i}$, noncommutative
Schatten--H\"older applied to
$A_{i}=E_{i}e^{-\ell_{i}H_{0}}$ with exponents $p_{i}=\beta/\ell_{i}$ gives the
fixed-history bound
\begin{equation}
\bigl|\operatorname{Tr}(A_{1}B_{1}\cdots A_{k}B_{k})\bigr|
\le\Bigl(\prod_{i}\|B_{i}\|\Bigr)
\exp\Bigl[-\frac D2\sum_{i}\ell_{i}\,\#\{x:\ a_{i,x}=h\}\Bigr].
\label{eq:schatten}
\end{equation}
The point of evaluating the heat-trace factors at $\beta$ rather than at the
individual $\ell_{i}$ is that $\operatorname{Tr}(Q_{x}^{h}e^{-\ell h_{x}})$
diverges as $\ell\to0^{+}$, so a bound organized by short intervals cannot be
written down, while the global Schatten exponents make the total high-site time
$-\frac D2\sum_{i}\ell_{i}\#\{x:a_{i,x}=h\}$ appear instead. Inserting a canonical
electronic projector of norm one preserves \eqref{eq:schatten}.

\subsection{The missing lemma, and why Miyao's theorem does not apply}

The ordering statement one would like is that the staggered density has a
positive two-point floor,
\begin{equation}
\bigl\langle M_{L}^{2}\bigr\rangle_{\beta,L}\ge m_{*}^{2},\qquad
M_{L}=N^{-1}\sum_{x}\eta_{x}s_{x},\qquad s_{x}=n_{x}-1,
\label{eq:ML}
\end{equation}
uniformly in even $L$, in the canonical half-filled sector, for some explicit
$m_{*}>0$ in a temperature window. \emph{This is open for the pure spinful
dispersionless Holstein model.} Note also that a one-point staggered density
vanishes in a symmetric finite-volume state and is not this order parameter, and
that at $t=0$ the canonical random-pair mixture has
$\langle M_{L}^{2}\rangle=1/(N-1)$, so the cold degeneracy alone carries no
order.

What \eqref{eq:LS}--\eqref{eq:schatten} provide is a local operator series and a
bound on a fixed history. What is missing is the normalized combinatorial
transfer, in the following precise form.

\begin{quote}\itshape
$W=0$ decorated-contour lemma. A volume-uniform, normalized decorated-contour
expansion for the full transformed Hamiltonian at the two separated scales $D$ and
$\Delta$, uniform in $L$ and $\beta$, supplying: (i) a space-time contour
representation in the canonical sector, or an explicitly justified conversion to
it, with the two checkerboards as bulk phases and the unrestricted high sector as
a single label; (ii) absolutely convergent connected vacuum-cluster weights whose
homogeneous part reproduces the order-$t^{2}$ coefficient and its bulk corrections
without double counting the transformed second-order term; (iii) after subtracting
the homogeneous bulk potential, a boundary-decoration bound
$\sum_{C:\,\widetilde C\cap\operatorname{core}\gamma\ne\varnothing}
|\Phi^{T}(C;\gamma)|\le\epsilon_{b}|\operatorname{core}\gamma|$, with the sum
taken over actual integrated connected-cluster weights and retaining the
difference of compatibility indicators; (iv) an anchored polymer majorant of the
form $\sup_{a}\sum_{\gamma:a\in\operatorname{core}\gamma}|w(\gamma)|e^{2|\operatorname{core}\gamma|}\le\frac18$,
with explicit lower bounds on the phase magnetization and control of
periodic and canonical finite-volume interfaces; (v) constants $c_{1},c_{2}$ and a
declared frequency interval such that the estimates hold whenever $t/U\le c_{1}$
and $k_{B}T\le c_{2}\Delta$.
\end{quote}

Three comments make the difficulty concrete. First, global particle-hole symmetry
does not help at an interface. It interchanges the two homogeneous checkerboards
and equates their bulk free energies; it does not equate the loop partition
function in a domain-wall background with that in a homogeneous background,
because a cluster touching the interface has different compatibility constraints,
and flipping only the interior changes the hopping that crosses its boundary. The
actual decorated-cluster formula of the Koteck\'y--Ueltschi type \cite{ku99}
retains a term
proportional to $\mathbf{1}_{C\sim\Gamma}-\mathbf{1}_{\{\text{background near }C
\text{ is a ground configuration},\,C\sim\text{that background}\}}$, which
vanishes when the cluster box misses the contour core and not when the two pure
phases are symmetry related. Second, a norm-small cold remainder does not control
the order-$t$ hopping inside the high block, repeated high excursions, or the
off-diagonal parts of the second-order term. Third, the closest rigorous bosonic
contour theorem, Miyao \cite{miyao24} with its correction \cite{miyao24corr},
requires a bare nearest-neighbour repulsion $W$ with $0<\varepsilon<W$, so it
excludes $W=0$ by hypothesis and does not cover \eqref{eq:cdwH}. Earlier
infrared-bound results in this family likewise assume a positive direct
interaction and three or more dimensions. Finite-volume singlet theorems are not
charge-order results.

This is a barrier of exactly the same shape as the one in
Section~\ref{sec:openproblem}: a strong-coupling effective theory plus a
volume-uniform local remainder does not by itself give a statement about the full
Gibbs state, and the missing object is again a normalized sum over decorated
histories with a controlled interface. The two open lemmas are not the same lemma,
and this note proves no implication between them. It records that the pairing
ceiling and the charge-order floor both stop at the same place, and that neither
can be reached from the effective Hamiltonian side by operator estimates alone.

\section{The remaining statement: a charged-skeleton influence bound}
\label{sec:openproblem}

\subsection{Why exactly one lemma is left}

Collecting the obstructions gives a short list of what a proof of \eqref{eq:lift}
cannot use. It cannot rotate the interactions away at operator level
(Section~\ref{sec:routes}\ref{sub:routea}), because the full generator is
resonant and the virtual cloud is of weight $\Theta(x^{2})$ at phonon number of
order $X$. It cannot dominate the influence functional pointwise by a Gaussian
(Section~\ref{sec:routes}\ref{sub:routeb}), because the only available
replacement discards the factor that carries $\Omega$. It cannot appeal to
reflection positivity (Section~\ref{sec:routes}\ref{sub:routec}) or to
log-concavity of the phonon marginal (Section~\ref{sec:routes}\ref{sub:routed}),
since both fail on two-site examples. It cannot factorize the cutoff-free kernel
conditionally in time (Section~\ref{sec:kernel}), because an explicit legal
history does better than the best free directed block by a margin that grows with
$U/\Omega$. And it cannot ignore the hot sector as rare, by
Section~\ref{sec:vtail}.

What survives is localized. The facing-plateau lemma of
Section~\ref{sec:kernel} supplies an explicit local cost of $3UW/22$ on each
facing plateau, with $W=\min(L_{0},1/\Omega)$ and $L_{0}$ the plateau length;
the Duhamel remainder of
Section~\ref{sec:vtail} supplies a convergent bound on short virtual histories;
and the directed kernel supplies the correct exponential scale
$e^{-X\coth(\beta\Omega/4)}$ for a free directed block. The missing object has to
combine these three into one statement about histories, and its precise form is
the following.

\subsection{The lemma}

A \emph{charged skeleton} is a space-time set
$S=\{(e_{i},u_{i})\}_{i=1}^{N}$ of pair-transfer events, on distinct oriented
bonds of a fixed space-time cylinder, with consecutive events separated in
imaginary time by at least a dwell scale $L_{0}>0$ and in space by a distance
monotone in $N$. Call the skeleton
\emph{persistent} if in addition every pair of consecutive events is connected by
a charged segment, that is, by a history whose charge on the intermediate region
is nonzero for the whole dwell interval.

\begin{quote}\itshape
Charged-skeleton influence bound. There are constants $\eta>0$ and $K<\infty$,
independent of volume and of $\beta\Omega\in(0,b]$, such that for every persistent
skeleton $S$ with $N$ events, the sum over all neutral decorations of the
screened weight satisfies
\[
\Bigl|\ \sum_{\text{neutral decorations of }S}\ e^{-\Phi_{\beta}[q]}\,
(\text{fermionic sign})\Bigr|
\ \le\ e^{-\eta N+K}\,\varrho\text{-weight}(S),
\]
where the right-hand side is normalized by the true partition function, the
constant $\eta$ may be chosen at least $c\,\frac{3}{22}UW$ for a bounded-overlap
constant $c>0$, so that the plateau costs are allocated to distinct events, and the sum
over short nonpersistent histories is bounded by a remainder of Duhamel type that
is uniform in volume.
\end{quote}

Four features are essential and none can be dropped. The sum is over signed
neutral histories and is normalized by $\varrho$, not by a product of free traces:
Section~\ref{sec:kernel} shows that the product of free directed weights is not an
upper bound for a legal history, so a normalized signed sum is the weakest object
that can still be true. The bound is on an \emph{influence}, not on a weight:
Section~\ref{sec:vtail} shows the weight is algebraically large. The cost is
allocated along the skeleton with bounded overlap, because the plateau gain is
local and does not add by itself. And the remainder is uniform in volume, which is
exactly what the local route of Theorem~\ref{thm:vtailT} achieves on its side and
what a volume-extensive Duhamel estimate does not.

\subsection{What is not claimed}

No implication is proved here in either direction between the bound just stated
and the lifted ceiling \eqref{eq:lift}. This note does not show that the bound is
true, does not show that it is sufficient, and does not show that it is necessary.
It shows that every route examined either fails on an explicit witness or reduces
to a statement of this shape. The distortion-valley route of
Section~\ref{sec:valley} ends at a locality step of the same class, in which
barrier crossings on adjacent bonds couple through the fermion determinant at
distance; that step is not shown here to follow from the lemma above. Whether a
persistent skeleton can be defined so that
the allocation of plateau costs is both possible and tight remains the central
open question, and it is a question about the geometry of long-lived pair
transport, not about the size of the phonon cloud.

\section{Scope of verification, data, and disclosures}
\label{sec:scope}

\paragraph{What is proved and what is computed.} Every theorem in this note is
accompanied by a proof. Three kinds of auxiliary evidence appear and are labelled
where they arise: exact algebraic identities, verified symbolically or in exact
rational arithmetic; interval-enclosed numerical evaluations of finite matrix
elements and finite sums; and converged finite-Fock exact diagonalization of small
systems, which is numerical evidence and not a proof of any infinite-dimensional
statement. No exact diagonalization result is used as a hypothesis of any theorem.

\paragraph{Replay scripts.} The arithmetic underlying
Theorem~\ref{thm:dir}, Proposition~\ref{prop:witness},
Theorem~\ref{thm:family}, Proposition~\ref{prop:loop} and
Theorem~\ref{thm:plateau} is reproduced by one script, the witnesses of
Section~\ref{sec:routes} by a second, the channel coefficients
\eqref{eq:wLF}--\eqref{eq:wlimit} together with the dimer diagonalization of
Section~\ref{sec:vtail} by a third, the coefficients, identities, window
constants and trace witnesses of Section~\ref{sec:cdw} by a fourth and an
independent fifth, and the Born-Oppenheimer surfaces, the splitting thresholds,
the oscillator-basis diagonalization of the dimer and the locality witness of
Section~\ref{sec:valley} by a sixth. The route-witness script, the charge-order
script and the distortion-valley script in addition reject deliberately broken
variants, nine, twelve and six respectively, each of which must fail. All six are
included in the supplement with instructions, and each exits with status zero
when run standalone. The scripts are checks of finite arithmetic; they are not
certificates for any thermodynamic statement.

\paragraph{Relation to the released results.} The cold-projected ceiling of
\cite{largeu} is used as an input and is not re-derived. The obstruction families
of Section~\ref{sec:routes} are stated for the estimates examined, not for all
possible estimates of that name. Nothing in this note changes the status of any
released result.

\paragraph{Use of software and language models.} AI tools were used for drafting,
calculation and checking under the author's direction and review. All theorems
were checked by hand against the proofs, and all numerical claims are reproducible
from the scripts described above.

This preprint has not been peer reviewed.

\paragraph{Data availability.} The manuscript source, the bibliography and the
six replay scripts are archived with this note. Released program results are
cited by DOI: \cite{largeu,strong,barriers2,barriers1,linear5}.

\bibliographystyle{plain}
\bibliography{refs}

\end{document}
